% !TEX root = ../../mixed-sklyanin-cubics.tex
% Current result: 2026-11-22 / v60.
\section{多中心の全滑らか領域、中央切断と可換比較}
\label{mct:sec:multicentral}

一中心で構成した実比較を、複数の共有中心一次変数へ拡張する。
まず因子別 PBW 関係と指定中心三次の全許容係数を分類する。
中央切断の判別を記述し、固定した滑らかな境界から実随伴を用いて
残りの全横断係数の comonad 比較を構成する。
横断二項 Bockstein の全階数から、一般係数の一節点性と判別因子の正規化も得る。
さらに厳密な中心作用から実 blow-up 圏を作り、
残余線束圏の smoothness と、一中心の係数式を偏極した可換候補の滑らかさを同値とする。
これにより全ての smooth な因子別対象に一般中央 flag が存在し、
任意個の中心で実際の可換三次への全 coherent Ku 比較を得る。
偏極候補と逐次比較の標的が同じであることまでは結論しない。
最後に二つの実制限の交換を調べ、軸上のデータだけでは失われる混合情報を明示する。

% !TEX root = ../../mixed-sklyanin-cubics.tex
% Current result: 2026-11-22 / v61.
\subsection{複数の共有中心変数に対する PBW 表示と全中心持上げ}
\label{mcp:subsec:pbw}

中心変数を増やす際、一次・二次係数を別々の一変数族に分けるだけでは、
中心三次の混合係数は決まらない。
ここでは既存の普遍中心恒等式を係数環の準同型で移し、混合項も含む全表示を得る。
また、複数の中心変数の線型結合が中心的であることから、各係数の中心性を
個別に結論することはしない。中心イデアルの随伴次数を使って全持上げを分類する。

本小節の代数的主張では $k\ge1$、$m\ge2$ とし、
滑らかな標準点曲線を持つ $S_i=S(A_i,B_i,C_i)$ と
$a_i\in\C^\times$ を固定する。
記号 $I_{\mathrm n},I_2$ はそれぞれ非二捩れ・二捩れ因子の集合とし、
$i\in I_2$ では $S_i=S(1,1,\tau_i)$ と規格化する。
各因子の生成元と標準中心三次 $\Omega_{i,0}$ も固定し、
\[
 B_0=\bigotimes_iS_i,\qquad
 w_0=\sum_i a_i\Omega_{i,0},\qquad D=B_0/(w_0)
\]
と置く。係数環 $R$ は任意の可換 $\C$-代数でよく、被約性は仮定しない。
共有中心変数を $\mathbf u=(u_1,\ldots,u_m)$、$\deg u_s=1$ とし、
係数 $R$ の内部次数は零とする。
以下の $R[\mathbf u]_d$ は中心変数について斉次次数 $d$ の部分を表す。

\begin{lemma}[普遍中心恒等式の多中心への代入]
\label{mcp:lem:polarization}
$i\in I_{\mathrm n}$ に対し、命題\ref{eco:prop:universal-lift}の語代表を固定し、
その普遍中心を補助 homogenizing 変数 $h$ で
\[
 \Omega_{i,0}+hC_{i,2}(P)+h^2C_{i,1}(P,L)
\]
と書く。$P$ は二次 cyclic potential、$L$ は一次 cyclic potential である。
$P(\mathbf u)$ の係数を $R[\mathbf u]_1$、
$L(\mathbf u)$ の係数を $R[\mathbf u]_2$ に取ると、
\begin{equation}\label{mcp:eq:substituted-center}
 \Omega_i(\mathbf u)=\Omega_{i,0}
       +C_{i,2}\bigl(P(\mathbf u)\bigr)
       +C_{i,1}\bigl(P(\mathbf u),L(\mathbf u)\bigr)
\end{equation}
は potential $\Phi_{i,3}+P(\mathbf u)+L(\mathbf u)$ の代数における
中心三次である。この代数は $R[\mathbf u]$ 上に固定した $S_i$ の PBW 基底を持つ。
この構成と基底は、非被約環への変更を含む任意の係数変更に両立する。
\end{lemma}

\begin{proof}
まず普遍係数環 $T_i=\C[P,L]$ 上で命題\ref{eco:prop:universal-lift}を用いる。
そこでの代数は $T_i[h]$ 上自由で、中心性はこの代数内の多項式恒等式である。
引用命題の証明では多項式環 $T_i$ の被約性を使ってこの恒等式を確定している。
以下ではその証明済みの恒等式を底変更し、$R$ の閉点での消滅から零性を推論しない。
$h=1$ として、$T_i$ 上自由な filtered 代数と中心元
\[
       \Omega_{i,0}+C_{i,2}(P)+C_{i,1}(P,L)
\]
を得る。次に $T_i\to R[\mathbf u]$ を、各 potential 係数に
指定した一次式・二次式を代入する準同型とする。
代数の表示、自由基底、中心恒等式をこの準同型で tensor するだけで、所要の主張を得る。
同命題では $C_{i,2}$ は $P$ に一次で active 次数二、
$C_{i,1}$ は係数重み二で active 次数一である。
従って式\eqref{mcp:eq:substituted-center}は総内部次数三である。
この段階では方向ごとの特殊化から中心性を推論しておらず、
$R$ の被約性や $u_s$ の可逆性は必要ない。
\end{proof}

混合係数を明示するため $C_{i,1}(P,L)=Q_i(P)+J_i(L)$ と書く。
$Q_i$ は二次写像、$J_i$ は線型写像であり、
$H_i(P,P')=(Q_i(P+P')-Q_i(P)-Q_i(P'))/2$ をその対称双線型偏極とする。
$m=2$ の場合
\[
 P(u,v)=uP_u+vP_v,\qquad
 L(u,v)=u^2L_{20}+uvL_{11}+v^2L_{02}
\]
なら式\eqref{mcp:eq:substituted-center}は
\begin{equation}\label{mcp:eq:two-variable-polarization}
 \begin{aligned}
 \Omega_i(u,v)={}&\Omega_{i,0}
   +uC_{i,2}(P_u)+vC_{i,2}(P_v)\\
 &+u^2\bigl(Q_i(P_u)+J_i(L_{20})\bigr)\\
 &+uv\bigl(2H_i(P_u,P_v)+J_i(L_{11})\bigr)\\
 &+v^2\bigl(Q_i(P_v)+J_i(L_{02})\bigr).
 \end{aligned}
\end{equation}
特に $uv$ 項には二つの一次 potential の偏極項が必要である。
任意個の変数でも、各 $u_su_t$ の係数を同じ双線型式で与える。

\begin{theorem}[指定境界を持つ全因子別 PBW・中心三次の正規形]
\label{mcp:thm:normal-form}
各因子の三関係に、その因子の生成元だけを使った
\begin{equation}\label{mcp:eq:general-relations}
       r_{i,j}+\sum_l M_{i,jl}(\mathbf u)x_{i,l}
                         -\alpha_{i,j}(\mathbf u)
\end{equation}
を許す。ただし $M_{i,jl}\in R[\mathbf u]_1$、
$\alpha_{i,j}\in R[\mathbf u]_2$ とし、異なる因子は可換とする。
この表示は、active 濾過の随伴代数が
$B_0\otimes_{\C}R[\mathbf u]$ となる意味で PBW であると仮定する。

この代数に $w\bmod(\mathbf u)=w_0$ を満たす中心三次 $w$ が存在するための
必要十分条件と、その全ての中心三次は、平行移動
$x_{i,j}=X_{i,j}+s_{i,j}(\mathbf u)$、$s_{i,j}\in R[\mathbf u]_1$ の後に、
次のように記述される。
\begin{enumerate}
 \item $i\in I_{\mathrm n}$ では $M_i$ は任意の対角零対称行列である。
 その三つの対角外成分を $p_i(\mathbf u),q_i(\mathbf u),r_i(\mathbf u)$ とし、
 \[
 P_i(\mathbf u)=p_i(\mathbf u)x_iy_i+q_i(\mathbf u)y_iz_i
                  +r_i(\mathbf u)z_ix_i,\qquad
 L_i(\mathbf u)=-\sum_j\alpha_{i,j}(\mathbf u)x_{i,j}
 \]
 と置く。中心元 $\Omega_i(\mathbf u)$ は
 式\eqref{mcp:eq:substituted-center}で与える。
 \item $i\in I_2$ では $M_i=0$ であり、定数補正 $\alpha_{i,j}(\mathbf u)$ は任意である。
 三つの平方はこの全係数上で中心的であり、中心三次として
 \[
          \Omega_{\tau_i}+\sum_j\beta_{i,j}(\mathbf u)x_{i,j}^2,
             \qquad \beta_{i,j}(\mathbf u)\in R[\mathbf u]_1
 \]
 を取る。
 \item 全体の中心三次は正確に
 \begin{equation}\label{mcp:eq:all-centers}
  \begin{aligned}
  w={}&\sum_{i\in I_{\mathrm n}}a_i\Omega_i(\mathbf u)\\
    &+\sum_{i\in I_2}a_i
      \left(\Omega_{\tau_i}+\sum_j\beta_{i,j}(\mathbf u)x_{i,j}^2\right)
    +c(\mathbf u),\qquad c(\mathbf u)\in R[\mathbf u]_3.
  \end{aligned}
 \end{equation}
 で尽くされる。
\end{enumerate}
対角零とする平行移動は一意であり、その後の上記全係数も一意である。
従って、固定した因子・生成元・中心境界を保つ表示をこの平行移動で割った族は、
\begin{equation}\label{mcp:eq:parameter-count}
          3km+3k\binom{m+1}{2}+\binom{m+2}{3}
\end{equation}
個の自由な係数を持つ affine 空間で表される。
$m=2$ ではその次元は $15k+4$ である。
この分類は任意の可換 $\C$-代数 $R$ に自然である。
\end{theorem}

\begin{proof}
補題\ref{ebc:lem:translation-gauge}の Koszul overlap を
$R[\mathbf u]$ 上で計算する。
PBW 条件は $M_i(\mathbf u)=M_i(\mathbf u)^{\mathsf t}$ を与える。
各 $C_i\ne0$ なので、その証明で使った
$R_i\cap\bigwedge^2V_i=0$ の線型恒等式は任意の係数環への変更でも成立する。
逆に対称行列は二次 cyclic potential の微分であり、
普遍 lower-potential 代数の PBW 基底を底変更すれば全係数で PBW となる。

平行移動は係数ごとに同補題の式\eqref{ebc:eq:translation}を与える。
その対角は $2C_i s_{i,j}(\mathbf u)$ だから、
$s_{i,j}=-M_{i,jj}/(2C_i)$ が対角を消す唯一の一次式である。
定数補正はこの代入で二次式に変わるが、依然として任意に許した範囲に入る。

次に、二捩れ因子の必要条件を中心イデアル $I=(u_1,\ldots,u_m)$ の一次層で検出する。
$M_i(\mathbf u)=\sum_s u_s M_i^{(s)}$ と書く。
中心性 $[x_{i,j},w]=0$ を $I^2$ を法として取り、
各 $u_s$ の係数と、他因子の active 次数が全て零となる純 $i$ ブロックへ射影する。
その式は
\[
       a_i\delta_{i,j}(M_i^{(s)})+[x_{i,j},Q_{i,s}]=0
                  \quad\text{in }S_{i,3}\otimes R
\]
となる。$w$ に含まれる混合ブロックの二次補正は、この純ブロック成分へ寄与しない。
定理\ref{eco:thm:clifford-obstruction}の平行移動による三つの線型な正規係数へ移し、
同定理の汎関数\eqref{eco:eq:three-functionals}を適用すると、
各 $M_i^{(s)}$ は平行移動の像に属さなければならない。
これは $\C$ 上の線型恒等式に $R$ を tensor したものであり、
閉点だけの判定ではない。$a_i\tau_i\ne0$ は $R$ 内でも単元なので、
非被約な係数は新しい解を与えない。
対角零の正規化後は各 $M_i^{(s)}=0$ である。

十分性について、非二捩れ因子には補題\ref{mcp:lem:polarization}を用いる。
二捩れ因子については、命題\ref{eco:prop:clifford-center-freedom}を
まず補助変数 $h=1$ とし、$\alpha_{i,j}\mapsto\alpha_{i,j}(\mathbf u)$ と底変更する。
その普遍 Clifford 表示は、平方と体積元の中心性を全 $R[\mathbf u]$ 上で与える。
従って式\eqref{mcp:eq:all-centers}の全ての元が中心的である。

最後に、これ以外の中心三次がないことを確認する。
正規化後の周囲代数を $B$ とし、$\beta=c=0$ で構成した中心元を $w^{\mathrm{can}}$ と書く。
$z=w-w^{\mathrm{can}}$ は総次数三の中心元で $z\in IB$ である。
PBW 基底により、中心イデアルについての随伴代数は
\begin{equation}\label{mcp:eq:i-associated}
       \operatorname{gr}_I B
             =B_0\otimes_{\C}R[\bar u_1,\ldots,\bar u_m]
\end{equation}
である。これは中心変数の次数による随伴代数であり、active 濾過とは区別する。
$z\in I^qB$ の初項は中心性により
\[
    R[\bar{\mathbf u}]_q\otimes_{\C}Z(B_0)_{3-q}
\]
の中に入る。
有限次数の中心条件は、生成元との交換子による有限個の $\C$-線型写像の核なので、
この核の計算は任意の $R$ への tensor と両立する。

体上の tensor 積の中心は各中心の tensor 積であり、
全因子で $Z(S_i)_1=0$、非二捩れでは $Z(S_i)_2=0$、
二捩れでは $Z(S_i)_2=\langle x_i^2,y_i^2,z_i^2\rangle$ である。
従って $q=1$ の初項は、二捩れ因子の平方に中心一次式を掛けたものだけである。
それらの平方は変形後の $B$ でも正確に中心的なので、その持上げを引けば
$z\in I^2B$ となる。
次に $Z(B_0)_1=0$ より $q=2$ の初項は零で、$z\in I^3B$ となる。
総次数三なので、残りは $R[\mathbf u]_3$ だけである。
ここで $I^4B$ の総次数三部分は零である。
以上により全ての自由度が式\eqref{mcp:eq:all-centers}に含まれる。
各次数層の基底独立性は、表示した係数の一意性も与える。
因子あたり三つの一次式と三つの二次式、最後の一つの三次式の係数を数えると
式\eqref{mcp:eq:parameter-count}を得る。
\end{proof}

\begin{proposition}[商の自由性と任意順序の中心正則列]
\label{mcp:prop:free-quotient}
定理\ref{mcp:thm:normal-form}の全係数族について、$A=B/(w)$ と置く。
$B,A$ は、次数付き加群としてそれぞれ
\begin{equation}\label{mcp:eq:free-bases}
       B\simeq R[\mathbf u]\otimes_{\C}B_0,\qquad
       A\simeq R[\mathbf u]\otimes_{\C}D
\end{equation}
という自由基底を持つ。この表示は代数同型を主張するものではない。
$w$ は $B$ の左右の非零因子であり、中心列 $u_1,\ldots,u_m$ は
$B,A$ の双方で任意の順序の正則列となる。
さらに $w,u_1,\ldots,u_m$ は $B$ 上で任意の順序の正則列である。
中心変数の可逆な $R$-線型変更を行っても同じ結論が成立する。
従って任意の部分集合 $J\subset\{1,\ldots,m\}$ に対し
\[
 A/(u_j:j\in J)
\]
は残りの中心多項式環上自由であり、全ての係数変更とこれらの商が可換する。
特に $A/(\mathbf u)=D\otimes_{\C}R$ である。

普遍係数環を $T_m$ とすれば各内部次数は有限自由 $T_m$-加群で、
その階数 Hilbert 級数は
\begin{equation}\label{mcp:eq:hilbert}
       H_B(z)=\frac1{(1-z)^{3k+m}},\qquad
       H_A(z)=\frac{1-z^3}{(1-z)^{3k+m}}
\end{equation}
となる。
\end{proposition}

\begin{proof}
$B$ の自由基底は普遍 PBW 基底の底変更から得られる。
まず $w_0$ が $B_0$ 上正則であることを確認する。
各 $S_i$ は $\C[\Omega_{i,0}]$ 上自由である。
実際、$S_i/(\Omega_{i,0})$ の斉次基底を持ち上げると、
次数に関する帰納法と $\Omega_{i,0}$ の正則性によりその基底を得る。
従って $B_0$ は $\C[\Omega_{1,0},\ldots,\Omega_{k,0}]$ 上自由であり、
非零線型式 $w_0=\sum_i a_i\Omega_{i,0}$ は左右正則である。
同じことは任意の $R$ を tensor しても成立する。

active 次数による濾過では
\[
      \operatorname{gr}_{\mathrm{act}}B=B_0\otimes R[\mathbf u],
         \qquad \operatorname{in}_{\mathrm{act}}(w)=w_0
\]
である。初項が非零因子なので $w$ も左右の非零因子である。
また $w$ の中心性により、その二側イデアルの全ての元は $wb$ の形であり、
$\operatorname{in}(wb)=w_0\operatorname{in}(b)$ となる。
従って商の濾過は厳密で、
\begin{equation}\label{mcp:eq:active-quotient}
           \operatorname{gr}_{\mathrm{act}}A
                 =D\otimes_{\C}R[\mathbf u]
\end{equation}
を得る。

$D$ の固定した斉次 $\C$-基底を語の線型結合で代表し、同じ語を $A$ へ持ち上げる。
式\eqref{mcp:eq:active-quotient}による active 次数の帰納法から、これらは
$R[\mathbf u]$ 上で $A$ を生成する。
有限な線型関係があれば、その最高 active 次数を取ると $D$ の基底間の関係になる。
係数は零なので、次数を下げて全ての係数が零であると分かる。
これが式\eqref{mcp:eq:free-bases}の第二式である。
この基底とその証明は普遍係数上で固定でき、その後の任意の底変更にも保たれる。

自由な $R[\mathbf u]$-加群上で、変数を任意の順序で消していく列は正則である。
また、任意の中心変数の部分集合を先に零とした $B$ 上でも、
$w$ の最高 active 部分は同じ $w_0$ である。
従って $w$ はその商でも正則であり、これで $w$ を途中に置く全ての順序も扱える。
可逆な中心座標変更は $R[\mathbf u]$ の自己同型であり、
係数一次式・二次式・三次式の形を保つので、同じ自由基底の証明を適用できる。
固定された graded 基底の個数から式\eqref{mcp:eq:hilbert}を得る。
\end{proof}

\begin{remark}[被覆の範囲と検算]
\label{mcp:rem:scope}
定理\ref{mcp:thm:normal-form}が分類するのは、因子間の可換性を保ち、
各因子の関係にその因子の active 一次項と中心二次項を加える指定表示である。
固定 tensor 境界と厳密な中心変数を持つ全 PBW 表示については、
定理\ref{ifp:thm:factorization}がこの因子別条件を強制する。
最高次の二次境界自体の変形や、任意の内在的非可換三次の表示被覆は含まない。
中心方向の多項式恒等式と全順序の正則性も、それだけでは残余圏の比較核や
複数の制限間の全 coherence を与えない。

\path{computations/verify_multicentral_pbw.py} は二中心変数について、
平行移動後の混合二次係数、既存の $S(1,-1,-1)$ 普遍中心公式の重み付き代入と
偏極による $uv$ 成分を厳密に確認する。
任意個の中心変数、任意の非被約係数環、全次数の自由基底と全中心持上げの分類は
上の証明によるものであり、この有限計算の確認範囲と区別する。
\end{remark}

% !TEX root = ../../mixed-sklyanin-cubics.tex
% Current result: 2026-11-22 / v59.
\subsection{中央切断の判別形式と比較を始められる開集合}
\label{mcpn:subsec:pencil}

以下は $k\ge2$、$N=3k$、中心変数二つ $U=(u,v)$ の場合である。
前節で一度固定した境界比較と Jacobian 切断から、
式\eqref{euk:eq:coefficients}の多項式 $Q,L,C$ を固定する。
定理\ref{mcp:thm:normal-form}の係数を
\[
 \mathbf t(U)\in\C^N\otimes\C[u,v]_1,\quad
 \alpha(U)\in\C^N\otimes\C[u,v]_2,\quad c(U)\in\C[u,v]_3
\]
と書く。これらを重み斉次式へ代入して得る三次式を
\begin{equation}\label{mcpn:eq:polarized-candidate}
 \begin{split}
 W_{\mathrm{pol}}(\mathbf x,U)
  ={}&G(\mathbf x)+Q(\mathbf t(U))(\mathbf x)\\
    &+L(\mathbf t(U),\alpha(U))(\mathbf x)
      +C(\mathbf t(U),\alpha(U),c(U)),\\
 X_{\mathrm{pol}}={}&V(W_{\mathrm{pol}})\subset\mathbb P^{N+1}
 \end{split}
\end{equation}
と定める。$Q,L,C$ の active 次数はそれぞれ二、一、零であり、
中心変数に関する次数は一、二、三である。
本小節はこの\emph{偏極した可換候補}の幾何を扱う。
多項式を偏極しただけで二中心の Morita 核が得られたとはしない。

一中心の判別多項式を系\ref{euk:cor:discriminant}の $D$ とし、
$\mu=2^N=8^k$ と置く。

\begin{proposition}[中央切断を判定する次数 $3\cdot8^k$ の形式]
\label{mcpn:prop:ray-discriminant}
\[
 \Delta(u,v)=D(\mathbf t(u,v),\alpha(u,v),c(u,v))
\]
は次数 $3\mu$ の二変数斉次形式である。恒等零となる特殊係数も許す。
$\ell=[a:b]\in\mathbb P^1$ に対し
\[
 A_\ell=A/(bu-av)
\]
と書くと、
\[
 \Delta(a,b)\ne0
 \quad\Longleftrightarrow\quad
 A_\ell\text{ の残余線束圏が smooth}
 \quad\Longleftrightarrow\quad
 V\bigl(W_{\mathrm{pol}}(\mathbf x,as,bs)\bigr)
       \subset\mathbb P^N\text{ が滑らか}
\]
である。従って全二中心係数空間の開集合
\[
 \mathcal U_{\mathrm{cut}}
   =\{\Delta\not\equiv0\}
\]
は、少なくとも一つの滑らかな中央切断を持つ係数を正確に表す。
その補集合は $\Delta$ の全 $3\mu+1$ 個の係数の消滅で定義される。
\end{proposition}

\begin{proof}
一中心側で active 変数を $\mathbf x\mapsto z\mathbf x$ と変えると、
係数の重み $(1,2,3)$ により臨界値も定数項も $z^3$ 倍される。
従って $z\ne0$ について
\[
 D(z\mathbf t,z^2\alpha,z^3c)=z^{3\mu}D(\mathbf t,\alpha,c).
\]
これは全係数の多項式恒等式なので $z=0$ にも成立する。
代入した $\Delta$ の斉次性が従う。
正規形の中心持上げは普遍多項式の底変更であるから、
$(u,v)=(as,bs)$ の代入は、一中心族の係数
$(\mathbf t(a,b),\alpha(a,b),c(a,b))$ と正確に一致する。
系\ref{euk:cor:discriminant}が表示した同値を与える。
非零斉次形式は $\mathbb P^1$ 全体では消えないので、最後の主張も従う。
\end{proof}

\begin{theorem}[一次形式を固定した一般係数の中央 pencil]
\label{mcpn:thm:lefschetz-pencil}
任意の $\mathbf t(U)$ を固定する。
二次形式 $\alpha(U)$ と三次形式 $c(U)$ の係数空間には非空 Zariski 開集合があり、
そこで $\Delta$ は正確に $3\mu$ 個の相異なる零点を持つ。
各零点の可換切断は唯一の通常二重点を持つ。
この開集合では $X_{\mathrm{pol}}$ は滑らかで、
\[
 \operatorname{Bl}_{Z}X_{\mathrm{pol}}\longrightarrow\mathbb P^1,
 \qquad Z=V(G)\subset\{u=v=0\},
\]
は $3\mu$ 個の一節点ファイバーを持つ Lefschetz pencil の解消である。
\end{theorem}

\begin{proof}
係数空間を
$\mathcal P=(\C^N\otimes H^0(\mathbb P^1,\OO(2)))
\times H^0(\mathbb P^1,\OO(3))$ とする。
各 $\ell\in\mathbb P^1$ で、値の評価
$(\alpha(U),c(U))\mapsto(\alpha(\ell),c(\ell))$ は全射である。
さらに $c$ の値と一次 jet は独立に指定できる。

固定した $\mathbf t(\ell)$ に対する $D=0$ は、
系\ref{euk:cor:normalization}により既約被約超曲面である。
その特異部分は $(\alpha,c)$ 空間で余次元二以上である。
値の評価の全射性から、$\mathbb P^1\times\mathcal P$ 上で
この特異部分へ入る incidence も余次元二以上となる。
従ってその $\mathcal P$ への像は真の閉部分集合である。
閉性には $\mathbb P^1$ の proper 性を使う。

この部分を除いた $D=0$ の点では、対応する臨界点が唯一で、
Hessian が非退化である。臨界代数をその長さ一の因子と残りに分けると、
norm の微分は
\[
                   \frac{\partial D}{\partial c}\ne0
\]
となる。ここでは $c\mapsto C$ の係数 $\kappa\ne0$ を使う。
$\Delta(\ell)=0$ を固定したまま $c$ の一次 jet を変えれば、
$d\Delta(\ell)$ を任意に変えられる。
従って $\Delta=d\Delta=0$ の incidence も、この領域では余次元二である。
前の特異部分を含めた閉集合の全次元は高々 $\dim\mathcal P-1$ なので、
その像を除いた非空開集合では全ての零点が単純となる。
次数 $3\mu$ から零点数を得る。

$Z$ 上では $G$ の偏微分が同時に消えないため、
$X_{\mathrm{pol}}$ と全ての切断は滑らかである。
$Z$ の外では中心方向の局所座標を $z$ とし、三次式を
$F(\mathbf x,z)=0$ と書く。切断が滑らかなら total space も滑らかである。
切断の唯一の Morse 点 $p$ では、上の長さ一因子による norm の局所表示から
\[
 d\Delta/dz\ne0\quad\Longrightarrow\quad
                       \partial F/\partial z(p)\ne0.
\]
従ってそこでも total space は滑らかである。
陰関数定理で $z$ を解くと、その Hessian は非退化なので
射影の臨界点は通常の Lefschetz 型となる。

$(u,v)$ は $X_{\mathrm{pol}}$ 上の正則列であり、
その零集合は滑らかな余次元二の $Z$ である。
従って pencil の graph はこの ideal の blow-up に一致する。
方向 $[a:b]$ のファイバーは
$V(W_{\mathrm{pol}}(\mathbf x,as,bs))$ 自身である。
各切断内では $Z$ は Cartier 因子なので余分な blow-up は生じない。
\end{proof}

次数 $3\cdot2^N$ 自体は、滑らかな $N$ 次元三次の双対次数という既知の数である
\cite[Section~7.1, Example~7.4]{CMOTevelevDuality}。
上の証明で確認したのは、$G$ と任意の $\mathbf t(U)$ を固定し、
制限された $Q$ の族の中でも、$\alpha(U),c(U)$ の変化だけで
この一般位置が得られることである。

\begin{example}[六 active 変数における $192$ 個の単純臨界方向]
\label{mcpn:ex:fermat-pencil}
$G=\sum_{j=0}^{5}x_j^3$、$b_j=2^j$ として
\[
 W=\sum_{j=0}^{5}x_j^3-3u^2\sum_{j=0}^{5}b_j^2x_j+v^3
\]
を取る。$u\ne0$ の臨界点は $x_j=\epsilon_jb_ju$、$\epsilon_j\in\{\pm1\}$ であり、
中央切断の判別形式は
\[
 \Delta(u,v)=\prod_{\epsilon\in\{\pm1\}^6}
          \left(v^3-2\sum_{j=0}^{5}\epsilon_j8^j\,u^3\right)
\]
である。符号付き和は互いに異なり、どれも零ではない。
従って $64$ 個の三次因子は互いに素で、各因子の三根は相異なる。
全 $192$ 方向が単純である。Hessian は
$\operatorname{diag}(6\epsilon_jb_ju)$ なので全臨界点が Morse であり、
$u=0$ のファイバーも滑らかである。
この有限の恒等式と jet の評価は
\texttt{computations/verify\_multicentral\_pencil.py}で再現する。
これは一般の比較係数の数値計算や、二中心核の検算ではない。
\end{example}

\begin{remark}[切断の判別と全対象の比較の区別]
\label{mcpn:rem:scope}
$\mathcal U_{\mathrm{cut}}$ 上で選んだ良い切断は、次の逐次比較の出発点になる。
そこで構成する標的と $X_{\mathrm{pol}}$ が同じであることは未証明である。
特に $X_{\mathrm{pol}}$ の smoothness を、そのまま二中心の非可換対象の
smoothness と同一視しない。
$\Delta\equiv0$ の部分に滑らかな非可換対象が存在しないことも、ここでは証明していない。
一般の非可換 Bertini を適用するには追加の確認が要る。
例えば\cite[Proposition~2.1]{CMORSVBertini}は可換有限型環上の有限性を仮定し、
同論文の Introduction は smooth dg 圏の層への拡張を証明していないと明記する。
一般の多中心 PBW 族にその有限中心仮定を無条件で加えることはしない。
\end{remark}

% !TEX root = ../../mixed-sklyanin-cubics.tex
% Current result: 2026-11-22 / v59.
\subsection{一つの滑らかな中央切断からの二中心比較}
\label{mci:subsec:iterated-comparison}

中心変数が二つある場合でも、その一つの中央切断が滑らかな残余圏を持てば、
一中心の比較を境界としてもう一度用いることができる。
この逐次構成は、中心平面の各 ray 上で選んだ核の貼合せではない。
一つの実境界核を固定した後、残る全横断係数の多項式環上で
実随伴の comonad を比較する。
以下では $k\ge2$、$N=3k$ とし、因子・生成元・非零重みを固定する。

\subsubsection{多中心の周囲解消と中央超平面制限}

\begin{proposition}[共有中心差を用いる相対線束模型]
\label{mci:prop:ambient-line-model}
中心変数の個数を $m\ge1$ とする。
$m=1$ では式\eqref{emc:eq:family}、$m\ge2$ では定理\ref{mcp:thm:normal-form}の
普遍係数環、または一部の係数を複素定数に固定した多項式係数環を $R$ とし、
その族を $A_m=B_m/(w)$ と書く。
$B_m,A_m$ は左右 noetherian で、各内部次数は有限自由 $R$-加群である。
$B_m$ は相対 Koszul であり、各複素ファイバーは
global dimension と Gorenstein parameter がともに $N+m$ の AS-regular 代数となる。
相対 augmentation 双対は
\begin{equation}\label{mci:eq:augmentation-duals}
 \begin{split}
 R\underline{\operatorname{Hom}}_{B_m}(R,B_m)
     &\simeq R(N+m)[-N-m],\\
 R\underline{\operatorname{Hom}}_{A_m}(R,A_m)
     &\simeq R(N+m-3)[-(N+m-1)]
 \end{split}
\end{equation}
であり、Nakayama twist はない。

全 twist の線束圏 $\mathcal C_m$ と、$a_m=N+m-3$ 本の例外線束を除いた
残余線束圏 $\mathcal K_m$ は、$R$-proper な有限生成 dg 模型を持つ。
複素数係数の任意の非零中央一次式を座標 $u_m$ として取り、
$A_{m-1}=A_m/(u_m)$ とすると、
制限は実 dg 関手
\[
 F_m:\mathcal K_m\longrightarrow\mathcal K_{m-1}
\]
とその左右随伴を与える。その twist と unshifted cotwist は
\begin{equation}\label{mci:eq:spherical-periods}
 T_{F_m}\simeq\varrho_{m-1}^2,\qquad
 C_{F_m}\simeq\varrho_m^{-1}[1],\qquad
 T_{F_m}^3\simeq[4],\qquad C_{F_m}^3\simeq[1]
\end{equation}
であり、標準随伴比較射も同型である。また
\begin{equation}\label{mci:eq:generation}
 \mathcal K_m=\operatorname{thick}(F_m^R\mathcal K_{m-1}),\qquad
 \mathcal K_{m-1}=\operatorname{thick}(F_m\mathcal K_m)
\end{equation}
となる。これらの実関手、随伴、周期と有限生成の構成は任意の導来底変更に保たれる。
同じ主張は、原点だけに affine 臨界点を持つ固定三次式を中央境界とする
可換三次族にも成立する。
\end{proposition}

\begin{proof}
一中心では命題\ref{emc:prop:line-model}、多中心では命題\ref{mcp:prop:free-quotient}が
$R[\mathbf u]$ 上の PBW 基底、商の基底、
$w$ と中心列の全順序の正則性を与える。
中心変数の線型変更後も同じ主張を適用できる。
$B_m/(\mathbf u)=B_0\otimes R$ は、主稿の補題
\paperref{main}{mbnp:lem:noetherian}と中央多項式拡大により左右 noetherian である。
正次数の中心元に沿う初項の持上げは内部次数ごとに有限回で止まるので、
中心変数を順に戻すと $B_m$ とその商 $A_m$ の noetherian 性を得る。

相対自己双対を明示する。$Q_0=R[u_1,\ldots,u_m]$ と置く。
各因子は $Q_0$ 上の cyclic potential の三微分で定まり、
cyclic Hessian による長さ三の自由双加群複体を持つ。
二捩れ因子の中心平方係数は商の中心元だけを変え、周囲の解消には現れない。
この複体の左右因子交換を伴う転置は $(3)[-3]$ 自己双対を与える。
active 濾過の随伴は固定した $S_i$ の Koszul 双加群解消であり、
各総内部次数の濾過が有限なので元の複体も exact である。
用いている自己双対複体は
\cite[Theorem~3.6, Lemma~3.7 and proof]{NFMBeTa}の potential 解消であり、
斉次標準因子の untwisted 性には\cite[Corollary~2.9]{MoriSmith}を用いる。
$Q_0$ 上の自由性から、各因子を tensor した長さ $N$ の複体も exact となる。

ここで
\[
 C=B_m\otimes_R B_m^{\mathrm{op}},\qquad
 D=B_m\otimes_{Q_0}B_m^{\mathrm{op}}
\]
とする。$C$ は $R[\mathbf u_L,\mathbf u_R]$ 上自由であり、
$z_s=u_{s,L}-u_{s,R}$ $(1\le s\le m)$ は中心正則列である。
従って $D=C/(z_1,\ldots,z_m)$ は長さ $m$ の Koszul 解消を持ち、
\[
 R\underline{\operatorname{Hom}}_C(D,C)\simeq D(m)[-m].
\]
相対 potential 解消の双対と change of rings により
\[
 R\underline{\operatorname{Hom}}_C(B_m,C)
          \simeq B_m(N+m)[-N-m]
\]
を得る。共有する各中心差を一度ずつ加えており、因子ごとに重複させていない。
各 $D$-自由項を上の有限 $C$-自由解消で置き換えると、
$B_m$ の $C$-perfect 性も得られる。

体値の内部次数別 bar 複体に active 濾過を入れると、随伴代数
$B_0\otimes\C[\mathbf u]$ の Tor は対角次数に集中する。
従って $B_m$ は Koszul であり、Hilbert 級数 $(1-z)^{-N-m}$ から
Koszul 空間の階数は $\binom{N+m}{j}$ となる。
普遍係数上では有限次数の行列の階数が全複素点で一定なので、
Koszul 空間は有限射影であり、その解消は底変更と両立する。
自己双対を片側へ augmentation し、正則中心三次の二項解消を用いると
式\eqref{mci:eq:augmentation-duals}を得る。
体値での有限双加群解消は global dimension の上界 $N+m$ を与え、
最上 Ext の非零性が等号を与える。

線束模型への入力も次数を追って確認する。
射影次元は $n_m=N+m-2=a_m+1$ であり、線束間の Ext は
\[
 \operatorname{Ext}^j(\OO(i),\OO(l))=
 \begin{cases}
  (A_m)_{l-i},&j=0,\\
  ((A_m)_{i-l-a_m})^\vee,&j=n_m,\\
  0,&\text{その他}
 \end{cases}
\]
となる。局所 cohomology と有限線束生成の証明は
命題\ref{emc:prop:line-model}の自己双対解消によるものを用い、
ここでは $N+m$ 本の連続する線束が全 twist を生成する。
各 Hom は $R$-perfect であり、有限例外列による mutation も底変更に両立する。

随伴射の計算が中心変数の個数に依存しないことを記す。
$a=a_m$、$i_*:\mathcal C_{m-1}\to\mathcal C_m$ を inflation とし、
残余包含の左右随伴を $p_L,p_R$ とする。
$i_*\OO(j)$ は $[\OO(j-1)\xrightarrow{u_m}\OO(j)]$ で解消される。
従って
\[
 F_m^R=p_Ri_*,\qquad F_m^L=p_Li_*(1)[-1].
\]
$M\in\mathcal K_m$ の制限に対して、$0\le r\le a-2$ の Hom 消滅は
$\operatorname{RHom}(\OO(r),M)$ と
$\operatorname{RHom}(\OO(r+1),M)$ の消滅から従う。
単位三角の残余射影は
\[
 C_{F_m}M=p_R(M(-1))[1]=\varrho_m^{-1}M[1]
\]
である。Serre 関手 $(-a)[a+1]$ により $\OO(-1)\in\mathcal K_m^\perp$ であり、
inverse rotation は $R_{\OO(-1)}\circ(-1)$ となる。
Orlov window の実周期 $(3)\simeq[2]$ はこの mutation 模型にも保たれる。

$Y=i_*N(1)[-1]$、$N\in\mathcal K_{m-1}$ とすると、
$1\le r\le a-1$ に対する
$\operatorname{RHom}(\OO(r),Y)
=\operatorname{RHom}(\OO(r-1),N)[-1]$ は零である。
従って $Y\to p_LY$ の correction は $\OO$ の生成圏に入り、
$(-1)[1]$ の後で $p_R$ がこれを殺す。
中心元の二項解消とその双対は、標準射
$F_m^R\to C_{F_m}F_m^L$ を同型にする。
また、$-a\le r\le-2$ に対して
\[
 \operatorname{RHom}(i_*N,\OO(r))
        =\operatorname{RHom}(N,\OO(r+1))[-1]=0
\]
なので、$p_Ri_*N\to i_*N$ の correction は $\OO(-1)$ だけである。
余単位の octahedron は $N(-1)[2]$ の right mutation を与え、
$T_{F_m}=\varrho_{m-1}^{-1}[2]=\varrho_{m-1}^2$ となる。
\cite[Theorem~5.1]{CMOAnnoLogvinenkoSpherical}の判定により spherical 性を得る。

最後に、proper 性から十分高い shift への射は零になる。
余単位三角を反復した $N\to T_{F_m}^{3q}N=N[4q]$、
および単位の fiber を反復した $M[-2q]\to M$ の cone を用いると、
式\eqref{mci:eq:generation}を得る。
この証明は全 coherent 圏の proper 性や source の smoothness を用いない。
有限生成対象についてこの有限構成を一度選べば、そのまま導来底変更できる。
可換側では周囲を多項式環として同じ解消と mutation を用いる。
\end{proof}

\subsubsection{実随伴からの rotation と一般変数数の arity}

\begin{lemma}[固定した滑らかな切断への周期付き rotation 移送]
\label{mci:lem:rotation-transfer}
一中心代数 $A_1$ と滑らかな可換三次 $Y_1$ の間に、
定理\ref{euk:thm:universal-kernel}の実 Morita 比較
$E_1:\mathcal K_1\simeq\Ku(Y_1)$ を選ぶ。
この比較から
\[
             \eta:E_1\varrho_1\xrightarrow{\sim}\varrho_{Y_1}E_1
\]
を構成でき、その三回の合成が両側の指定周期 $\varrho^3\simeq[2]$ を保つように取れる。
因子とこの複素係数点を固定すれば、$E_1,\eta$ は任意の横断係数環へ定数族として移せる。
\end{lemma}

\begin{proof}
最初の中央制限を $F_u$ と書く。
上記の実比較は、固定された二側の境界同値 $E_0$ とともに
$E_0F_u\simeq F_u^{\mathrm{cl}}E_1$ および実随伴の単位・余単位を保つ。
右随伴を取り、単位三角を移すと
\[
 E_1C_{F_u}\simeq C_{F_u^{\mathrm{cl}}}E_1
\]
を得る。式\eqref{mci:eq:spherical-periods}からこれは
$E_1\varrho_1^{-1}\simeq\varrho_{Y_1}^{-1}E_1$ を与え、逆を取って所要の
rotation 比較 $\eta_0$ を得る。

二つの指定三乗周期と $\eta_0$ の三重合成の差は、可逆核 $E_1[2]$ の自己同型である。
BFK の対角公式\cite[Theorem~5.39]{NFBFKKernels}から
$\operatorname{HH}^0(\Ku(Y_1))=\C$、$\operatorname{HH}^{<0}(\Ku(Y_1))=0$ となる。
実際、恒等 sector の次数零は Jacobian 環の次数零であり、
非恒等 sector がこの次数へ寄与する条件は変数の個数が零であることになる。
負次数も両 sector で消える。
従って周期の差は一つの $b\in\C^\times$ であり、
$\eta_0$ を $d$ 倍すると差は $d^3b$ となる。
$d^3=b^{-1}$ と取れば周期が一致する。
二つの周期比較射はこれで $H^0$ において一致する。
その同定 path の $i$ 次 homotopy 群は
$H^{-i-1}\operatorname{RHom}(E_1[2],E_1[2])$ $(i\ge0)$ であり、全て零である。
従って同定の空間は可縮で、周期関係の合成に必要な高次 homotopy にも
追加の選択障害はない。実 dg 核上の周期付き比較として使える。

この scalar 補正後に、元の $u$ 作用の標識も同時に保つとは結論しない。
ここで必要なのは固定境界の圏、rotation、および三次周期の比較である。
全てを複素数体上で一度選んだ後、その核と射を係数環へ tensor する。
因子や切断を動かす底で三乗根を一斉に選ぶ議論ではない。
\end{proof}

\begin{lemma}[変数数を三の倍数に限らない比較群]
\label{mci:lem:smooth-cubic-arity}
$M\ge6$ とし、$F\in\C[y_1,\ldots,y_M]_3$ の射影零点集合 $Y$ を滑らかとする。
$V_d=\operatorname{Jac}(F)_d$、$\mathcal B=\Ku(Y)$ と書く。
任意の係数環 $R$ と $\eta\in V_2\otimes R$ に対して
\[
 K_\eta=\operatorname{fib}(\mathbf1\xrightarrow{\eta}\varrho^2),\qquad
 J=\varrho^2[-1]
\]
と置くと、$r\ge2$、$i\ge0$ について
\begin{equation}\label{mci:eq:arity}
 H^{2-r-i}\operatorname{RHom}(K_\eta,J^r)=
 \begin{cases}
  V_1\otimes R,&(r,i)=(2,0),\\
  R,&(r,i)=(3,0),\\
  0,&\text{その他}
 \end{cases}
\end{equation}
となる。この同定は $\eta$ と底変更に自然である。
固定 counit の正規化 comonad 塔の非空 fiber は離散であり、四項以降は可縮である。
また $H^{-j}\operatorname{RHom}(\mathbf1,\varrho^2)=0$ $(j\ge1)$ である。
\end{lemma}

\begin{proof}
BFK 公式\cite[Theorem~5.39 and proof, pp.~82--84]{NFBFKKernels}は
変数数の三による整除を仮定しない。
$H^a\operatorname{RHom}(\mathbf1,\varrho^b)$ の恒等 sector は
$a$ が偶数なら $V_{b+3a/2}$、奇数なら零である。
二つの非恒等 sector の必要十分条件は
\[
                  a\equiv M\pmod2,\qquad 2b+3a=M
\]
である。$U_r=\operatorname{RHom}(\mathbf1,J^r)$、
$W_r=\operatorname{RHom}(J,J^r)$ と書き、余単位三角を用いると、
$d=2-r-i$ に対して
\[
 H^{d-1}W_r\longrightarrow H^dU_r\longrightarrow
 H^d\operatorname{RHom}(K_\eta,J^r)\longrightarrow H^dW_r
\]
を得る。三つの既知項の恒等 sector の次数は順に
\[
 1-r-3i/2\quad(i\text{ 偶数}),\qquad
 3-r-3i/2\quad(i\text{ 偶数}),\qquad
 (5-2r-3i)/2\quad(i\text{ 奇数})
\]
であり、該当 parity 以外では零となる。
両端は負次数で、中項は式\eqref{mci:eq:arity}の二場合に限られる。
非恒等 sector の条件は順に
\[
 2-2r-3i=M,\qquad 6-2r-3i=M,\qquad 5-2r-3i=M
\]
であり、左辺が高々二なので不可能である。
両隣接群の消滅が $H^dU_r$ との自然な同型を与える。
最後の負次数消滅も $b=2,a=-j$ の同じ公式から従う。

$U_r,W_r$ は固定した複素数体上の複体の底変更であり、
関連 cohomology は自由である。非分裂の mapping complex 全体を一定とする必要はない。
正規化 unit-insertion cube の total fiber は $J^r$ であるから、
補題\ref{cmo:lem:normalized-comonad-uniqueness}で確認した opposite と Ind 完備化を通して
\cite[Theorem~3.3, Lemma~3.11]{CMOWillumsgaardAssociativity}を適用する。
その $i$ 次 homotopy 群が式\eqref{mci:eq:arity}であり、塔についての結論を得る。
\end{proof}

\subsubsection{固定した滑らかな ray と全横断係数}

二中心の正規形を $\mathbf U=(u,v)$ で書く。
非二捩れ因子の二次 potential 係数と、二捩れ因子の中心平方係数を
共通に $\mathbf t$ と書き、次の係数を用いる。
\begin{equation}\label{mci:eq:two-center-parameters}
 \begin{split}
 \mathbf t(u,v)&=u\mathbf t_u+v\mathbf t_v,\\
 \alpha(u,v)&=u^2\alpha_{20}+uv\alpha_{11}+v^2\alpha_{02},\\
 c(u,v)&=c_{30}u^3+c_{21}u^2v+c_{12}uv^2+c_{03}v^3.
 \end{split}
\end{equation}
まず $b=(\mathbf t_u,\alpha_{20},c_{30})$ を一つの複素係数として固定し、
$A_1=A_2/(v)$ に対応する定理\ref{euk:thm:universal-kernel}の標的
\begin{equation}\label{mci:eq:fixed-boundary}
 Y_1=V(F_1)\subset\mathbb P^N,\qquad
 F_1=G+uQ(\mathbf t_u)+u^2L(\mathbf t_u,\alpha_{20})
                  +u^3C(\mathbf t_u,\alpha_{20},c_{30})
\end{equation}
が滑らかであると仮定する。
これは系\ref{euk:cor:discriminant}の有限判別行列で判定できる明示的な条件である。
特に $\mathcal K_1$ は smooth proper であり、全 coherent 残余圏に一致する。

残る横断係数とその重みを
\begin{equation}\label{mci:eq:transverse-weights}
 \begin{gathered}
 \gamma_1=(\mathbf t_v,\alpha_{11},c_{21})\in\C^{2N+1},\qquad
 \gamma_2=(\alpha_{02},c_{12})\in\C^{N+1},\qquad z=c_{03},\\
 R=\C[\gamma_1,\gamma_2,z],\qquad
 \operatorname{wt}(\gamma_1)=1,\quad
 \operatorname{wt}(\gamma_2)=2,\quad\operatorname{wt}(z)=3
 \end{gathered}
\end{equation}
とする。この重みは $u$ を固定して $v$ を rescale するものに対応する。
$M=N+1$ とし、境界変数を $\mathbf y=(x_1,\ldots,x_N,u)$ と並べる。

\begin{theorem}[滑らかな一中央切断を持つ二中心族の実核]
\label{mci:thm:iterated-kernel}
式\eqref{mci:eq:fixed-boundary}が滑らかとなる $b$ を固定する。
全横断係数環 $R$ 上に
\begin{equation}\label{mci:eq:target}
 \begin{split}
 W_{\mathrm{it}}={}&F_1(\mathbf y)+v\widetilde Q(\gamma_1;\mathbf y)
       +v^2\widetilde L(\gamma_1,\gamma_2;\mathbf y)
       +v^3\widetilde C(\gamma_1,\gamma_2,z),\\
 \widetilde L={}&\mathsf B\gamma_2+\widetilde R_2(\gamma_1),\\
 \widetilde C={}&\kappa' z+\widetilde R_{1,1}(\gamma_1,\gamma_2)
                              +\widetilde R_3(\gamma_1)
 \end{split}
\end{equation}
という可換三次式が存在する。
ここで $\widetilde Q$ は $\gamma_1$ に線型、$\mathbf y$ に二次、
$\mathsf B:\C^{N+1}\to\operatorname{Jac}(F_1)_1$ は定数線型写像、
$\widetilde R_2$ は $\gamma_1$ に斉次二次で $\mathbf y$ に一次、
$\widetilde R_{1,1}$ は双線型、$\widetilde R_3$ は斉次三次、
$\kappa'\in\C^\times$ である。
$\mathsf B$ の可逆性は本定理では主張しない。

$X_{\mathrm{it}}=V(W_{\mathrm{it}})\subset\mathbb P^{N+1}$ とすると、
$A_2$ と $X_{\mathrm{it}}$ の残余線束圏の間に、全 $R$ 上の実 Morita 双加群と逆双加群が存在する。
これらは $v=0$ への実制限を、補題\ref{mci:lem:rotation-transfer}で選んだ $E_1$ を通して
単位・余単位とともに intertwine し、任意の導来底変更に保たれる。
各複素係数で
\[
 \mathcal K_2\text{ が smooth}\quad\Longleftrightarrow\quad
 X_{\mathrm{it}}\text{ が滑らか}
\]
である。この smooth 開上では核は相対積上 perfect となり、
各複素ファイバーで全 coherent 残余圏の同値
\[
          \Ku_{\mathrm{coh}}(A_2)\simeq\Ku(X_{\mathrm{it}})
\]
を得る。
\end{theorem}

\begin{proof}
境界 $A_1$ は横断係数に依存しない。
補題\ref{mci:lem:rotation-transfer}で周期付きの実核 $E_1$ を一度選び、$R$ へ tensor する。
命題\ref{mci:prop:ambient-line-model}は $v$ 制限の実随伴を全 $R$ 上で与える。
その twist 射を $E_1$ により
\[
       \eta\in\operatorname{Jac}(F_1)_2\otimes R
\]
へ移す。$v$ 一次の Bockstein は $A_2/(v^2)$ によって決まり、
式\eqref{mci:eq:transverse-weights}の重み一から $\gamma_1$ に線型となる。
固定線型切断
$\operatorname{Jac}(F_1)_2\to\C[\mathbf y]_2$ を選び、
実可換族 $F_1+v\widetilde Q$ の twist 射が $\eta$ となるように取る。
その可換 twist 射は $-[\widetilde Q]$ である。
これは境界 MF の微分 $d^2=F_1$ を持ち上げ、
$W=0$ 上で $d^2=-v(\widetilde Q+v\widetilde L+v^2\widetilde C)$ とする
補題\ref{euk:lem:commutative-reference}の二項解消そのものである。
同じ cohomology 類を表す二射の homotopy は、固定 mapping 複体の
重みを保つ $\C$-線型縮約を $R$ へ tensor して多項式的に選べる。
負次数 Hom の消滅により、その path の空間は可縮である。

次に、この実可換族を二項構造の参照点にする。
補題\ref{mci:lem:smooth-cubic-arity}の自然な同定により、二項差は
$\operatorname{Jac}(F_1)_1\otimes R$ の大域切断である。
可換側で $\widetilde L$ を加える作用は $-\operatorname{id}$ であり、
これは $v^2$ に対する法線二列解消の平方作用素から得られる。
そのため、この差を消す線型式 $\widetilde L$ を全 $R$ 上で選べる。
重み二から $\widetilde L=\mathsf B\gamma_2+\widetilde R_2(\gamma_1)$ となる。
ここでは $\mathsf B$ の全階数を仮定も結論もしていない。

二項を同定する path の空間は可縮である。
そこで実可換族 $F_1+v\widetilde Q+v^2\widetilde L$ を三項の参照点に取ると、
三項差は $R$ の大域切断となる。
可換側の $v^3$ 係数による作用は非零定数倍なので、これを消す
$\widetilde C$ を多項式的に選べる。
重み三の単項式は $z$、$\gamma_1\gamma_2$、$\gamma_1^3$ に限られるため、
式\eqref{mci:eq:target}の形を得る。

$\kappa'\ne0$ は $\gamma_1=\gamma_2=0$ で直接確認できる。
このとき source の周囲代数は固定した一中心の周囲代数の中央多項式拡大であり、
曲率は $w_1+zv^3$ である。
restriction--inflation の実解消の法線因子は
\[
 d_0=\begin{pmatrix}0&v\\zv^2&0\end{pmatrix},\qquad
 \theta=\begin{pmatrix}0&1\\-zv&0\end{pmatrix},\qquad
 h=\begin{pmatrix}0&0\\-z&0\end{pmatrix}
\]
となる。$\theta^2=d_0h+hd_0=-zv\operatorname{id}$、
$\theta h+h\theta=-z\operatorname{id}$ なので三項は $\pm z$ である。
この法線因子は実随伴核を解消したものであり、
境界の非零単位へ tensor すると
$\operatorname{Jac}(F_1)_0=\C$ の非零 class を与える。
固定された $E_1$ と周期比較はこの一次元空間の同型だから、
可換側との比 $\kappa'$ は非零である。

全係数上の二つの実構造の差を使ったので、閉点ごとの minimal model の貼合せは不要である。
四項以降は、補題\ref{mci:lem:smooth-cubic-arity}により非空比較 fiber が可縮となる。
両側に実随伴の全構造が存在するため、三項までの比較 path を全高次へ延ばせる。
これで全 homotopy-coherent comonad の実同値を得る。

境界の生成対象を $G_1$ とする。
式\eqref{mci:eq:generation}により $F_2^RG_1$ は source の生成対象であり、
その自己準同型代数は
\[
 \operatorname{REnd}(F_2^RG_1)
       \simeq\operatorname{RHom}(K_\eta G_1,G_1),\qquad
       g\star f=gK_\eta(f)\Delta_{G_1}
\]
という実 coKleisli 模型である。
可換側も同じ生成性を持ち、全 comonad 比較が積と全高次構造を保つので、
graph 双加群とその逆を得る。
生成対象上の制限 intertwiner と単位・余単位は、有限 cone と直和因子へ延びる。
全ての構成を $R$ 上で行ったため、任意の導来底変更にも保たれる。

可換 smooth 開では、その残余圏と source の有限線束模型は相対 smooth proper となる。
graph は $R$-perfect であり、両側の smooth 対角の有限解消を用いると
相対積上も perfect となる。
命題\ref{mci:prop:ambient-line-model}の noetherian AS-regular 周囲代数と
正則中心三次を用い、主稿の命題\paperref{main}{crs:prop:residual-smooth}と
定理\paperref{main}{crs:thm:coherent}によって各複素ファイバーの全 coherent 残余圏へ飽和する。
逆に、可換残余線束圏が smooth なら有限例外列を戻した
$\operatorname{Perf}(X_{\mathrm{it}})$ も smooth である。
同じ飽和判定は全 coherent 対象を perfect にするから、
特異 hypersurface 局所環の剰余体の無限射影次元により特異点は存在できない。
これで smoothness の必要十分条件を得る。
\end{proof}

\begin{corollary}[最後の中心三次係数に関する有限判別]
\label{mci:cor:finite-discriminant}
前定理の $b$ を固定し、$R_0=\C[\gamma_1,\gamma_2]$ と置く。
\[
 f=F_1+\widetilde Q+\widetilde L,\qquad
 \mathcal Q=R_0[\mathbf y]/(\partial f),\qquad
 D_{\mathrm{it}}=\det(\widetilde C\operatorname{id}_{\mathcal Q}+m_f)
\]
とすると、$\mathcal Q$ は階数
\[
                     \mu'=2^{N+1}=2\cdot8^k
\]
の有限自由 $R_0$-加群であり、$[f]=[(\widetilde Q+2\widetilde L)/3]$ となる。
$D_{\mathrm{it}}$ は $z=c_{03}$ に関して次数 $\mu'$、最高係数 $(\kappa')^{\mu'}$ を持ち、
\[
 D_{\mathrm{it}}\ne0\quad\Longleftrightarrow\quad
 X_{\mathrm{it}}\text{ が滑らか}
 \quad\Longleftrightarrow\quad\mathcal K_2\text{ が smooth}
\]
となる。従って任意の固定した全横断係数 $\gamma_1,\gamma_2$ に対し、
高々 $2\cdot8^k$ 個の $c_{03}$ を除いて全 coherent Ku 比較を得る。
悪い係数は重複度を込めれば正確に $\mu'$ 個である。
\end{corollary}

\begin{proof}
$Y_1$ の滑らかさにより $\partial F_1$ は原点だけを共通零点とする二次正則列である。
active 変数をここでは $\mathbf y$ 全体として filtered Koszul 複体を用いると、
\[
            \operatorname{gr}\mathcal Q=\operatorname{Jac}(F_1)\otimes R_0
\]
を得る。固定 Jacobian 基底の持上げが $R_0$-基底となり、
Hilbert 級数 $(1+t)^{N+1}$ から階数を得る。
Euler の式 $3F_1+2\widetilde Q+\widetilde L=0$ が表示した $[f]$ を与える。
境界 $v=0$ は滑らかなので、全特異点は $v=1$ の chart にある。
そこでの条件は $\partial f=0$ と $f+\widetilde C=0$ であり、
有限代数の単元条件から行列式による判定を得る。
$\widetilde C$ の $z$ 係数は非零定数 $\kappa'$ だから、
全ての $\gamma_1,\gamma_2$ の特殊化後も最高係数は消えない。
これは多重臨界点・衝突した臨界値を含む norm であり、
相異なる根や一般の通常二重点性をこの段階では主張しない。
\end{proof}

\begin{corollary}[少なくとも一本の滑らかな中央切断]
\label{mci:cor:smooth-ray}
定理\ref{mcp:thm:normal-form}の任意の複素係数二中心代数を考える。
中心平面の一次元部分空間 $\ell$ に沿う一中心商について、
系\ref{euk:cor:discriminant}の判別値が非零となる $\ell$ が少なくとも一つ存在すれば、
元の二中心代数の残余線束圏は、同じ次元のある可換三次の残余線束圏と実 Morita 同値である。
その残余線束圏が smooth な場合、標的は滑らかとなり全 coherent Ku 比較を得る。
\end{corollary}

\begin{proof}
その $\ell$ を $v=0$ とする中心変数の線型座標を一度選ぶ。
命題\ref{mcp:prop:free-quotient}と正規形はこの変更後にも成立する。
$u$ 方向の係数を固定すれば式\eqref{mci:eq:fixed-boundary}が滑らかであるから、
前定理の全横断係数族に元の代数が一つのファイバーとして含まれる。
このファイバーへ構成した核を底変更すればよい。
\end{proof}

\begin{corollary}[滑らかな中央 flag を仮定する反復]
\label{mci:cor:smooth-flag}
任意の $m\ge2$ について、定理\ref{mcp:thm:normal-form}の複素係数代数に
中央 flag
\[
 A_m\longrightarrow A_{m-1}\longrightarrow\cdots
       \longrightarrow A_1\longrightarrow D,
 \qquad A_{j-1}=A_j/(u_j)
\]
を固定する。$1\le j<m$ の各 $A_j$ の残余線束圏が smooth であると仮定する。
このとき $A_m$ の残余線束圏は、同じ次元のある可換三次の残余線束圏と
実 Morita 同値である。最上段も smooth なら全 coherent Ku 同値を得る。
さらに、固定した $A_{m-1}$ を保つ横断係数の多項式環上では相対的な実核を構成できる。
\end{corollary}

\begin{proof}
初段は定理\ref{euk:thm:universal-kernel}とその smoothness 判定による。
各段の実制限 intertwiner と随伴から次の境界の cotwist を移し、
補題\ref{mci:lem:rotation-transfer}の scalar 三乗根で指定周期を同期する。
$j$ 段目の可換境界の変数数は $N+j-1\ge7$ であるから、
補題\ref{mci:lem:smooth-cubic-arity}をそのまま適用できる。
$u_j$ の一乗、二乗、三乗に付く横断係数をそれぞれ重み一、二、三としてまとめれば、
実可換参照族の二次・一次・定数補正が各段で全ての比較群を実現する。
その前段が smooth という仮定と、命題\ref{mci:prop:ambient-line-model}の実随伴を用い、
前定理の実可換参照族・全 comonad・coKleisli の構成を反復する。
最上段の係数だけを動かす場合、境界核を固定してからその係数環へ tensor するので、
同じ構成は相対核を与える。
flag 自体が存在するという主張は、この条件付き系には含まれない。
\end{proof}

\begin{remark}[偏極候補・ray の貼合せとの区別]
\label{mci:rem:scope}
一中心の係数多項式に $\mathbf t(u,v),\alpha(u,v),c(u,v)$ を代入した
素朴な偏極候補を $W_{\mathrm{pol}}$ と書くと、
本小節で構成した $W_{\mathrm{it}}$ がそれに一致することは証明していない。
$\widetilde Q,\mathsf B,\kappa',\widetilde R_2,
\widetilde R_{1,1},\widetilde R_3$ は、固定した滑らかな切断とその比較核の選択に依存する。
特に $\kappa'$ へ一中心の指定値 $\kappa_k^2=27/28$ を流用しない。
横断の重み二写像 $\mathsf B$ の可逆性と、一般係数での臨界値の相異性は別の問題である。

本小節は、固定した一本の smooth な中央切断を境界として、
その補空間方向の全係数を一つの実随伴で扱った。
全ての中央 ray の比較を $\mathbb P^1$ 上へ降下すること、
吹上げの例外成分と元の残余圏を核で同時に識別すること、
二つの中心作用を coherent に balance することは、この証明の入力ではない。
因子や flag を動かす大域核、全標識の同時保存、
全ての中央切断が特異となる係数の可換比較、および全 intrinsic 対象の被覆は
引き続き本結論の外にある。
\end{remark}

% !TEX root = ../../mixed-sklyanin-cubics.tex
% Current result: 2026-11-22 / v60.
\subsection{横断二項 Bockstein の全階数}
\label{mtb:subsec:transverse-bockstein}

定理\ref{mci:thm:iterated-kernel}の境界係数 $b$ と、周期を同期した実境界核
$E_1$ を固定する。本小節では、その横断二項写像 $\mathsf B$ が同型であることを示す。
定数関係の $N$ 方向は零項の曲率で、中心式の $uv^2$ 方向は一項の微分で検出される。
後者は対象ごとの obstruction が零であるため、零項の評価だけでは検出できない。

以下で $A_1=B_1/(w_1)$、$\mathbf y=(x_1,\ldots,x_N,u)$、$M=N+1$ と置く。
命題\ref{mci:prop:ambient-line-model}により、$B_1$ は $M$ 個の生成元を持つ
quadratic Koszul AS-regular 代数で、$w_1$ は正則な中心三次である。
その二次双対を $E=B_1^!$ と書く。
$B_1$ の二次関係空間には各因子の $N$ 本の関係 $r_{ij}^{(b)}$ と、
因子間および $u$ との交換関係が基底をなす。
前者の双対を $q_{ij}\in E_2$、生成元 $u$ の双対を $\xi_u\in E_1$ とする。

\begin{lemma}[固定した一中心境界の Tate 次数]
\label{mtb:lem:tate-degrees}
$\kappa$ を $A_1/(A_1)_{\ge1}=\C$ の安定化とし、
全 $\kappa(a)$、$a\in\mathbb Z$ を対象とする full dg 部分圏を $\mathcal T_1$ とする。
中央周期作用素 $\vartheta$ に関して
\begin{equation}\label{mtb:eq:tate-module}
 \underline{\operatorname{Ext}}_{A_1}^*(\C,\C)
       =E\otimes\C[\vartheta],\qquad
 H^*(\mathcal T_1)=E\otimes\C[\vartheta,\vartheta^{-1}],
 \qquad \deg E_j=(j,-j),\quad\deg\vartheta=(2,-3)
\end{equation}
となる。後式は全内部次数付き射空間をまとめて書いたものである。
特に双次数 $(2,-2)$ の部分は $E_2$、$(1,-2)$ の部分は零である。
\end{lemma}

\begin{proof}
補題\ref{ebc:lem:full-tate}の Koszul--Shamash 構成を $B_1,w_1$ に適用する。
同補題のこの部分は、周囲の quadratic Koszul 性と中心三次の正則性だけを用いる。
第 $n$ homotopy は homological degree を $2n-1$、内部次数を $3n$ 増し、
Koszul 自由項の間での係数次数は $n+1>0$ である。
従って augmentation 後の微分は零となり、前式を得る。
Shamash の divided-power shift は自己交叉双加群からの中央 cohomology operator
を表す。中心的な derived localization
\cite[Theorem~5.3, Corollary~5.16]{NFMBCLLocalization}と
同補題の Tate-cut により後式を得る。
これは高次積を自明とする主張ではない。

$E_j\vartheta^m$、$m\in\mathbb Z$ の双次数が $(2,-2)$ となる条件は
$j+2m=2$、$j+3m=2$ で、解は $(j,m)=(2,0)$ だけである。
$(1,-2)$ の場合は $j+2m=1$、$j+3m=2$ となり、
唯一の整数解 $(j,m)=(-1,1)$ は $E_{-1}=0$ により現れない。
\end{proof}

\begin{proposition}[横断二項写像の可逆性]
\label{mtb:prop:transverse-iso}
定理\ref{mci:thm:iterated-kernel}の任意の滑らかな固定境界 $b$ に対し、
\begin{equation}\label{mtb:eq:transverse-iso}
 \mathsf B:\C^{N+1}_{(\alpha_{02},c_{12})}
         \xrightarrow{\ \sim\ }\operatorname{Jac}(F_1)_1
\end{equation}
は線型同型である。固定した因子、境界係数または周期の同期について、
同定理以外の一般性条件は必要ない。
\end{proposition}

\begin{proof}
まず $\gamma_1=0$、$z=0$ とし、$\epsilon=v^2$、$\deg\epsilon=2$ と置く。
$\alpha=\alpha_{02}$、$d=c_{12}$ と省略する。
多中心正規形の普遍恒等式から、この族は中心変数 $\epsilon$ に関する族
\begin{equation}\label{mtb:eq:epsilon-family}
 \mathcal A_{\alpha,d}=\mathcal B_\alpha/
       \bigl(w_1+\epsilon(\Lambda_b(\alpha)+du)\bigr),
 \qquad r_{ij}^{(b)}-\epsilon\alpha_{ij}=0
\end{equation}
からの底変更である。交換関係は変わらず、
$\Lambda_b(\alpha)\in (B_1)_1$ は $\alpha$ に線型である。
二捩れ因子の標準体積元にはこの補正が不要である。
式\eqref{mtb:eq:epsilon-family}は $\C[\epsilon]$ 上自由であり、
$\epsilon$ は中心正則元である。
これは補題\ref{mcp:lem:polarization}と
命題\ref{mcp:prop:free-quotient}の active 初項の証明を、
内部次数二の係数 $\epsilon$ に適用して得られる。

\emph{実自己交叉からの full class.}
$\epsilon$ の二項自由解消の双加群自己交叉は
\[
 A_1(-2)[1]\longrightarrow
 A_1\otimes_{\mathcal A_{\alpha,d}}^{\mathbf L}A_1
 \longrightarrow A_1\xrightarrow{\eta_{\alpha,d}}A_1(-2)[2]
\]
を与える。補題\ref{ebc:lem:full-tate}の bar 比較により、
$\eta_{\alpha,d}$ は $\mathcal A_{\alpha,d}/(\epsilon^2)$ の積の
一次変形 cocycle の class である。
自由解消の圏への移送には
Lowen\cite[Theorems~3.21, 4.15 and proof]{CMOLowenCharacteristic}の
full brace cochain を使い、singularity quotient、window と
中央周期の derived localization を通す。
従って以下の零項と一項は、実自然変換
\[
 \operatorname{id}_{\mathcal K_1}\longrightarrow
       \varrho_1^{-2}[2]\simeq\varrho_1
\]
を $\mathcal T_1$ へ制限した full cochain の成分である。
対象ごとの obstruction だけを自然変換と同一視していない。

\emph{定数関係の $N$ 方向.}
Koszul--Shamash 解消を式\eqref{mtb:eq:epsilon-family}へ持ち上げ、
微分の二乗を $\epsilon$ で割って augmentation すると、二次関係成分は
\begin{equation}\label{mtb:eq:curvature-evaluation}
            \delta_0(\kappa)=\pm\sum_{i,j}\alpha_{ij}q_{ij}\in E_2
\end{equation}
となる。符号は二拡張の共通の規約による。
中心三次の $\epsilon$ 係数は一次式なので、この評価では零になる。
従って $d$ および $\Lambda_b(\alpha)$ は
式\eqref{mtb:eq:curvature-evaluation}に混入しない。
補題\ref{mtb:lem:tate-degrees}により、この双次数には周期を含む別成分がない。
各 $q_{ij}$ は二次関係基底の双対として独立であり、
$E[\vartheta]$ に $\vartheta$-torsion がないため安定化後も消えない。
よって $\eta_{\alpha,d}=0$ なら $\alpha=0$ である。

\emph{中心一次式 $du$ の一項検出.}
以後 $\alpha=0$ とする。このとき周囲代数は $B_1[\epsilon]$ であり、
$\mathbf y\mapsto0$ は
$\mathcal A_{0,d}\to\C[\epsilon]$ という augmentation を与える。
従って $\kappa$ とその全 twist を相対自由解消で持ち上げられ、
full cochain の零項は零に取れる。

$K_\bullet$ を $B_1$ 上の Koszul 解消とする。
その第零項から第一項への $w_1$ の nullhomotopy は二次係数を持つ。
$u$ の作用は augmentation 上で零であり、中央作用なので
$K_\bullet$ 上で nullhomotopic である。
その同次 homotopy $h_u$ の第零項は $1\mapsto e_u$ と選べる。
従って $w_1+\epsilon du$ の最初の Shamash homotopy は
\[
                    s_1^{(\epsilon)}=s_1+\epsilon d h_u
\]
と取れる。高次 homotopy は通常の再帰で補う。
相対解消に $\operatorname{Hom}(-,\C[\epsilon])$ を適用すると、
$K_0\to K_1$ の追加係数 $\epsilon d e_u$ の双対が
\begin{equation}\label{mtb:eq:unary-u}
 \delta_1(\xi_u)=\pm d\vartheta,\qquad
 \delta_1(\xi_{ij})=0
\end{equation}
を与える。
実際、第 $n$ homotopy の係数次数は $n+1$ であり、
$\epsilon$ 係数を取った後は $n-1$ である。
$n\ge2$ では正次数なので augmentation 後に消え、
式\eqref{mtb:eq:unary-u}へ別の高次項が寄与しない。
この式は相対 Hom 複体の Bockstein が cohomology に誘導する写像である。
中央周期作用素の局所化は full cochain と両立するため、
同じ式を最小 Tate 模型で用いることができる。

\emph{全 homotopy に対する非消滅.}
最小模型を $m_1=0$ とし、次数二・内部次数 $-2$ の cochain
$\delta$ が全 Hochschild 境界 $[m,h]$ であると仮定する。
primitive の零項 $h_{0,a}$ は、各対象 $\kappa(a)$ について
双次数 $(1,-2)$ の射である。
補題\ref{mtb:lem:tate-degrees}によりこの射空間は零であり、
対象ごとに別々の primitive を許しても $h_{0,a}=0$ である。
もとの圏は uncurved で $m_0=0$ だから、境界の一項成分は
\[
                   [m,h]_1=[m_2,h_0]+[m_1,h_1]=0
\]
となる。任意の高 arity 成分と全ての $\vartheta$ の逆冪を許しても、この式は変わらない。
式\eqref{mtb:eq:unary-u}の $\vartheta$ は非零なので $d=0$ となる。
全 $\mathcal K_1$ 上の nullhomotopy はこの full 部分圏へ制限されるため、
これで $\eta_{\alpha,d}$ の $N+1$ 方向の単射性が従う。

\emph{横断二項との同定.}
上の class が定理\ref{mci:thm:iterated-kernel}の二項差であることは、
$\epsilon=v^2$ の実二列法線解消で確認できる。
持上げ微分の二乗を $\widetilde d^{\,2}=\epsilon\widetilde\theta$ と書くと、
$\epsilon$ の正則性から $[\widetilde d,\widetilde\theta]=0$ である。
従って法線の exterior 生成元 $e$ に対し
\[
 D=\widetilde d+v(\partial_e-\widetilde\theta e),\qquad
 \xi=\partial_e+\widetilde\theta e,\qquad
                         \xi^2=\widetilde\theta
\]
となる。従って二項 comonad class は同じ $\epsilon$-Bockstein である。
この計算は命題\ref{ebc:prop:constant-iso}と同じ実随伴の計算であり、
固定境界に既に $u$ があることによって変わらない。

$E_1$ と指定周期によりこの class を可換境界へ移すと、
補題\ref{mci:lem:smooth-cubic-arity}が
$\operatorname{Hom}(\operatorname{id},\varrho_1)=\operatorname{Jac}(F_1)_1$
を与える。右辺の次元は $M=N+1$ である。
可換参照族の二項係数の作用は $-\operatorname{id}$ なので、
この移送の後にその逆を施した写像が $\mathsf B$ である。
単射性と次元から式\eqref{mtb:eq:transverse-iso}が従う。
境界核や周期同期の選択変更は自然変換空間の同型を与えるため、
この全階数の結論に影響しない。

最後に、横断二項の重みは二であり、$\gamma_1$ の重みは一、
$\gamma_2$ の重みは二である。
従って $\gamma_2$ の線型係数に正次数の $\gamma_1$ は現れず、
$\gamma_1=0$ で検出した同じ $\mathsf B$ が全横断係数上で現れる。
\end{proof}

\begin{corollary}[横断係数の三角多項式同型]
\label{mtb:cor:triangular-coefficients}
定理\ref{mci:thm:iterated-kernel}の固定境界の下で、写像
\[
 (\gamma_1,\gamma_2,z)\longmapsto
       (\gamma_1,\widetilde L,\widetilde C)
\]
は $\gamma_1$ の係数環上の多項式同型である。
従って任意の固定 $\gamma_1$ について、$\widetilde L$ の全 $N+1$ 個の一次係数と
$\widetilde C$ を独立に指定できる。
\end{corollary}

\begin{proof}
$\mathsf B$ と $\kappa'$ はともに可逆なので、逆写像は
\[
 \begin{split}
 \gamma_2&=\mathsf B^{-1}
       \bigl(\widetilde L-\widetilde R_2(\gamma_1)\bigr),\\
 z&=(\kappa')^{-1}
       \bigl(\widetilde C-\widetilde R_{1,1}(\gamma_1,\gamma_2)
                         -\widetilde R_3(\gamma_1)\bigr)
 \end{split}
\]
で与えられる。
\end{proof}

\begin{corollary}[任意個の固定中心に対する横断同型]
\label{mtb:cor:many-centers}
$q\ge1$ とし、多中心正規形の固定した境界 $A_q=B_q/(w_q)$ を考える。
中心変数を $\mathbf u=(u_1,\ldots,u_q)$ とし、$M=N+q$ と置く。
その残余線束圏が smooth であり、滑らかな三次
$Y_q=V(F_q)\subset\mathbb P^{M-1}$ への、指定周期を同期した実境界核 $E_q$ を固定する。
この入力は、系\ref{mci:cor:smooth-flag}の滑らかな中央 flag から得られる。

新しい中心変数 $v$ を加え、境界係数を固定して全許容横断係数を動かす。
重み二の係数は
\[
 \gamma_2=(\alpha_{v^2},d_1,\ldots,d_q)\in\C^{N+q},
 \qquad v^2\sum_{j=1}^q d_j u_j
\]
であり、後式は中心三次の当該補正を表す。
このとき逐次比較の二項写像は線型同型
\[
                 \mathsf B_q:\C^{N+q}\xrightarrow{\sim}
                                  \operatorname{Jac}(F_q)_1
\]
である。全重み一係数を $\gamma_1$、$v^3$ の係数を $z$ とまとめると、
対応する $(\gamma_1,\gamma_2,z)\mapsto
(\gamma_1,\widetilde L,\widetilde C)$ も三角多項式同型となる。
\end{corollary}

\begin{proof}
命題\ref{mci:prop:ambient-line-model}により、$B_q$ は $N+q$ 個の次数一生成元を持つ
quadratic Koszul AS-regular 代数である。
従って $E=B_q^!$ として補題\ref{mtb:lem:tate-degrees}と同じ Tate 模型を持つ。
$\gamma_1=z=0$、$\epsilon=v^2$ とした変形は
\[
 r_{ij}^{(b)}-\epsilon\alpha_{v^2,ij}=0,\qquad
 w_q+\epsilon\left(\Lambda_b(\alpha_{v^2})
                                +\sum_{j=1}^q d_j u_j\right)=0
\]
となる。零項の評価は $\pm\sum\alpha_{v^2,ij}q_{ij}$ であり、
独立な $N$ 本の二次関係双対によって $\alpha_{v^2}=0$ を先に得る。
残る場合には相対 augmentation があり、各 $u_j$ の nullhomotopy を用いて
\[
                    \delta_1(\xi_{u_j})=\pm d_j\vartheta
                    \quad(1\le j\le q)
\]
を得る。高次 homotopy の係数次数と primitive $h_0$ の双次数の計算は
生成元数に依存しないので、全 $d_j$ が零でなければ full class は消えない。
自己交叉から実二項差への橋も同じ $\epsilon=v^2$ の二列解消である。

固定した $E_q$ と周期で移すと、補題\ref{mci:lem:smooth-cubic-arity}の
$\operatorname{Jac}(F_q)_1$ は次元 $M=N+q$ であり、単射性から全単射性を得る。
全横断係数での実核と三項係数の非零性には、
系\ref{mci:cor:smooth-flag}の証明中の、固定境界から一段進める構成を用いる。
同じ重み $(1,2,3)$ の議論により二項の $\gamma_2$ 係数は一定であり、
三項の $z$ 係数も非零定数である。
従って系\ref{mtb:cor:triangular-coefficients}と同じ逆写像を持つ。
ここでは、比較核をまだ持たない任意の smooth 境界に
$E_q$ が自動的に存在するとは仮定していない。
\end{proof}

\begin{remark}[この検出が与えない同定]
\label{mtb:rem:scope}
本小節は一つの滑らかな境界と実比較核を固定して、横断二項の全階数を証明した。
この同型の標準座標での行列を求めたわけではなく、
$W_{\mathrm{it}}=W_{\mathrm{pol}}$ や元の二中心作用を同時に保つ核を構成したわけでもない。
全 smooth 対象に良い中央切断と flag が存在することは、
定理\ref{mcb:thm:central-flags}で別に証明する。
その定理は本小節の全階数を入力としない。
\end{remark}

% !TEX root = ../../mixed-sklyanin-cubics.tex
% Current result: 2026-11-22 / v60.
\subsection{横断判別超曲面の正規化と単純な最後の中心値}
\label{mtd:subsec:transverse-discriminant}

系\ref{mci:cor:finite-discriminant}の有限判別を、
命題\ref{mtb:prop:transverse-iso}の全階数によって強める。
境界係数 $b$ と実比較核は引き続き固定し、$Y_1=V(F_1)$ は滑らかとする。
$N=3k$、$M=N+1$、$\mu'=2^M=2\cdot8^k$ と置く。
本小節の臨界点は $\mathbf y=(x_1,\ldots,x_N,u)$ の全 $M$ 変数について取る。

系\ref{mtb:cor:triangular-coefficients}の座標変更を
\begin{equation}\label{mtd:eq:theta}
 \Theta:(\gamma_1,\gamma_2,z)\longmapsto
      (\gamma_1,\ell,C),\qquad
 \ell=\operatorname{coeff}_{\mathbf y}\widetilde L,\quad
 C=\widetilde C
\end{equation}
と書く。$\gamma_1$ の係数空間を $P=\mathbb A^{2N+1}$ とし、
\[
 H_{\gamma_1}=F_1+\widetilde Q(\gamma_1),\qquad
 f_{\gamma_1,\ell}=H_{\gamma_1}+\ell\cdot\mathbf y
\]
と置く。独立な座標 $(\gamma_1,\ell,C)$ 上で
\[
 \begin{gathered}
 R=\C[\gamma_1,\ell],\qquad
 \widehat{\mathcal Q}=R[\mathbf y]/(\partial f_{\gamma_1,\ell}),\\
 \widehat D(\gamma_1,\ell,C)
       =\det(C\operatorname{id}_{\widehat{\mathcal Q}}+m_f)
 \end{gathered}
\]
と定めると、系\ref{mci:cor:finite-discriminant}の判別式は
$D_{\mathrm{it}}=\widehat D\circ\Theta$ である。
同系により $\widehat{\mathcal Q}$ は階数 $\mu'$ の自由 $R$-加群である。

以下の議論は、補題\ref{ldg:lem:fixed-quadratic-morsification}と
定理\ref{ldg:thm:normalization}の、$M=N+1$ 変数への適用である。
それらの証明に変数数の三による整除は不要であることを、必要な箇所で確認する。
用いる条件は、最高次の $\partial F_1$ が原点だけを共通零点とする
二次式の正則列であることと、基礎体が標数零であることである。

\begin{lemma}[任意の固定一次係数での臨界値分離]
\label{mtd:lem:separated-critical-values}
任意の $\gamma_1^0\in P(\C)$ に対し、
$\mathbb A^M_\ell$ の非空 Zariski 開集合 $U_{\gamma_1^0}$ が存在し、
その各点で $f_{\gamma_1^0,\ell}$ は正確に $\mu'$ 個の臨界点を持つ。
これらは全て Hessian が非退化であり、臨界値も互いに相異なる。
\end{lemma}

\begin{proof}
$H=H_{\gamma_1^0}$ と省略する。臨界スキーム
\[
 \mathscr C=
 \operatorname{Spec}\C[\ell,\mathbf y]/(\partial_jH+\ell_j)_{j=1}^M
       \longrightarrow\mathbb A^M_\ell
\]
は $\ell=-\nabla H(\mathbf y)$ により $\mathbb A^M_{\mathbf y}$ と同型である。
最高次部分は $\partial F_1$ なので、filtered Koszul 複体による
系\ref{mci:cor:finite-discriminant}の証明は、この射が階数 $\mu'$ の有限平坦射であることを与える。
整スキーム間の有限優勢射であり、標数零なので関数体拡大は分離的である。
従って非空開集合 $U_0\subset\mathbb A^M_\ell$ 上で有限 \`etale となる。
その $\mu'$ 個の点で、Jacobian 行列 $-\operatorname{Hess}H$ は可逆である。

$\mathscr C_0=\mathscr C\times_{\mathbb A^M_\ell}U_0$ とし、
相異なる順序付き臨界点対の有限 \`etale スキーム
\[
 T=(\mathscr C_0\times_{U_0}\mathscr C_0)\setminus\Delta
\]
を考える。各非空既約成分は $U_0$ を支配する。
二つの点を $\mathbf y,\mathbf y'$ と書くと、臨界値差
$h=f_\ell(\mathbf y)-f_\ell(\mathbf y')$ は
\begin{equation}\label{mtd:eq:critical-value-differential}
                  dh=\sum_{j=1}^M(y_j-y_j')\,d\ell_j
\end{equation}
を満たす。$\Omega^1_{T/\C}$ は $d\ell_j$ を基底に持つから、
$h$ がある成分で恒等零ならその成分で全ての $y_j-y_j'$ が零となり、
対角を除いたことに矛盾する。
従って各成分の $V(h)$ は真の閉集合で、その有限像も
$U_0$ 全体にはならない閉集合である。
有限個のこれらの像を除いて $U_{\gamma_1^0}$ を得る。
ここでは二次項 $\widetilde Q(\gamma_1^0)$ に一般性を課していない。
\end{proof}

\begin{theorem}[横断判別因子の既約性と affine 正規化]
\label{mtd:thm:normalization}
$D_{\mathrm{it}}\in\C[\gamma_1,\gamma_2,z]$ は既約であり、
任意の固定した $\gamma_1^0$ について
$D_{\mathrm{it}}(\gamma_1^0,\gamma_2,z)\in\C[\gamma_2,z]$ も既約である。
従ってそれぞれが定める判別因子は被約である。

全横断係数の判別超曲面の正規化は $P\times\mathbb A^M_{\mathbf y}$ であり、
正規化射は
\begin{equation}\label{mtd:eq:normalization}
 (\gamma_1,\mathbf y)\longmapsto
 \Theta^{-1}\!\left(
 \gamma_1,-\nabla H_{\gamma_1}(\mathbf y),
 \mathbf y\cdot\nabla H_{\gamma_1}(\mathbf y)-H_{\gamma_1}(\mathbf y)
 \right)
\end{equation}
で与えられる。各固定 $\gamma_1^0$ の判別超曲面についても、
この式の制限 $\mathbb A^M_{\mathbf y}\to V(D_{\mathrm{it}}(\gamma_1^0,-,-))$
が正規化である。

複素点として、判別超曲面の特異点は次のいずれかを満たす係数と正確に一致する。
\begin{enumerate}
 \item $f_{\gamma_1,\ell}$ が値 $-C$ を取る相異なる臨界点を二つ以上持つ。
 \item 値 $-C$ を取る臨界点の少なくとも一つで Hessian が退化する。
\end{enumerate}
この判定は全横断係数と各固定 $\gamma_1^0$ のどちらにも成立する。
\end{theorem}

\begin{proof}
まず座標 $(\gamma_1,\ell,C)$ で証明する。
普遍臨界スキームは
\[
 \widehat{\mathscr C}=\operatorname{Spec}\widehat{\mathcal Q}
       \simeq P\times\mathbb A^M_{\mathbf y}
\]
であり、整かつ正規である。
$C$ を $-f$ へ送る環準同型は有限射
\[
 \nu(\gamma_1,\mathbf y)=
 \left(\gamma_1,-\nabla H_{\gamma_1}(\mathbf y),
  \mathbf y\cdot\nabla H_{\gamma_1}(\mathbf y)-H_{\gamma_1}(\mathbf y)\right)
\]
を定める。有限性は $\widehat{\mathcal Q}$ の $R$ 上の有限性から従う。
有限臨界代数で $C+f$ が非可逆であることと $\widehat D=0$ は同値なので、
この射の像の台は $V(\widehat D)$ である。

$K=\operatorname{Frac}R$ とすると
$L=\widehat{\mathcal Q}\otimes_RK$ は次数 $\mu'$ の体拡大となる。
元 $-f\in L$ の最小多項式を $m(T)$ とすれば、
その乗法作用の特性多項式は
\[
                 \widehat D(T)=m(T)^{[L:K(-f)]}
\]
である。前補題により $\widehat D$ が $C$ に関して単純根だけを持つ特殊化が存在するので、
一般係数体上でもこの特性多項式は分離的である。
従って指数は一で、$L=K(-f)$ となる。
$\widehat D$ は $C$ に関して monic なので、Gauss の補題により $R[C]$ 上既約である。
任意の固定 $\gamma_1^0$ でも前補題が使えるため、
$R=\C[\ell]$ とした同じ証明で既約性を得る。
一般の norm 多項式に既約性を仮定しているのではなく、
ここでは臨界スキームの整性と臨界値分離の両方を用いている。

$L=K(-f)$ は $\nu$ が双有理であることを意味する。
その源は正規で、射は有限だから、$\nu$ は判別超曲面の正規化射である。
固定 $\gamma_1^0$ の場合も同様である。
臨界スキームから $(\gamma_1,\ell)$ への射の次数は $\mu'$ だが、
判別超曲面への正規化射の次数は一である。

特異点の判定も確認する。$A=\operatorname{Hess}_{\mathbf y}H_{\gamma_1}$ と置くと、
\[
          \ker d\nu=\{(0,\mathbf v):A\mathbf v=0\}.
\]
固定 $\gamma_1$ の微分も核が $\ker A$ である。
また $\nu^{-1}(\gamma_1,\ell,C)$ の点は、
値 $-C$ を取る臨界点と一対一に対応する。
判別超曲面が滑らかな点では正規化が局所同型なので、
逆像は一点で、その Hessian は非退化でなければならない。

逆に逆像が一点で Hessian が非退化なら、有限な局所環の包含
$A_0\subset B_0$ は剰余体 $\C$ を共有し、余接写像が全射である。
両極大イデアルを $\mathfrak m,\mathfrak n$ とすると、
$\mathfrak n=\mathfrak mB_0+\mathfrak n^2$ から
Nakayama の補題で $\mathfrak n=\mathfrak mB_0$ を得る。
さらに $B_0=A_0+\mathfrak mB_0$ なので、有限 $A_0$-加群 $B_0/A_0$ に
再び Nakayama を適用して $A_0=B_0$ となる。
源は正則だから判別超曲面もその点で滑らかである。
これは定理\ref{ldg:thm:normalization}の局所環の議論であり、$M$ の値に依存しない。
最後に多項式同型 $\Theta$ で戻すと、全ての結論と式\eqref{mtd:eq:normalization}を得る。
\end{proof}

\begin{corollary}[各固定一次係数に対する $2\cdot8^k$ 個の単純な悪い値]
\label{mtd:cor:simple-values}
任意の固定した $\gamma_1$ に対して、$\gamma_2$ の非空 Zariski 開集合が存在し、
そこで $D_{\mathrm{it}}(\gamma_1,\gamma_2,z)$ は $z$ に関して
正確に $2\cdot8^k$ 個の相異なる単根を持つ。
各根に対応する可換三次標的 $X_{\mathrm{it}}$ は唯一の通常二重点を持ち、
その他の $z$ では滑らかである。
従ってこの一般 $\gamma_2$ の下では、同じ $2\cdot8^k$ 個の値が
元の残余線束圏の非 smooth な係数の全体である。
\end{corollary}

\begin{proof}
固定した $\gamma_1$ の下で
$\gamma_2\mapsto\ell=\mathsf B\gamma_2+
\operatorname{coeff}_{\mathbf y}\widetilde R_2(\gamma_1)$ は affine 線型同型である。
補題\ref{mtd:lem:separated-critical-values}の開集合の逆像を取る。
そこでは $f_{\gamma_1,\ell}$ の $\mu'$ 個の臨界点が全て Morse 点であり、
臨界値も相異なる。
固定 $\gamma_2$ に対する $z\mapsto C$ は係数 $\kappa'\ne0$ の affine 同型なので、
各臨界値はちょうど一つの単純な悪い $z$ に対応する。

$v=0$ の境界 $Y_1$ は滑らかだから、$X_{\mathrm{it}}$ の特異点は全て $v=1$ にある。
そこでの方程式は $f_{\gamma_1,\ell}+C=0$ である。
各悪い値で零となる臨界値は一つだけなので特異点は唯一であり、
Hessian の非退化性と形式的 Morse 補題から完備局所環は
\[
             \C[[s_1,\ldots,s_M]]/(s_1^2+\cdots+s_M^2)
\]
となる。最後の非可換側の smoothness 判定は
定理\ref{mci:thm:iterated-kernel}の実 Morita 比較による。
\end{proof}

\begin{corollary}[一般の中央段での最後の判別次数]
\label{mtd:cor:many-centers}
$m\ge2$ とし、系\ref{mtb:cor:many-centers}の入力を
$q=m-1$ 個の中心を持つ固定境界に適用する。
すなわち境界は smooth で、滑らかな可換三次への周期付き実核を固定する。
新しい中心変数 $v$ の全重み一係数 $\gamma_1$ を任意に固定すると、
重み二係数 $\gamma_2$ の非空 Zariski 開集合上で、
最後の $v^3$ 係数には正確に
\[
                              2^{N+m-1}
\]
個の相異なる単純な悪い値がある。
各可換標的は唯一の通常二重点を持ち、それ以外の値では滑らかである。

固定した任意の $\gamma_1$ について、$(\gamma_2,z)$ の判別多項式は既約であり、
その零因子は被約、正規化は $\mathbb A^{N+m-1}$ である。
全横断係数についても判別多項式は既約であり、正規化は
重み一係数の affine 空間と $\mathbb A^{N+m-1}$ の積である。
\end{corollary}

\begin{proof}
$M=N+m-1$ とする。系\ref{mtb:cor:many-centers}により、
固定 $\gamma_1$ の下で重み二係数は全 $\ell\in\C^M$ を動き、
$z$ は全定数項 $C$ を動かす。
滑らかな境界式 $F_q$ の $\partial F_q$ は二次正則列なので、
臨界代数の階数は $2^M$ である。
補題\ref{mtd:lem:separated-critical-values}の証明は、
この有限平坦な gradient 写像と標数零だけを用いているためそのまま適用できる。
臨界値差の式\eqref{mtd:eq:critical-value-differential}も変わらない。
従って単純根と唯一の通常二重点を得る。
臨界スキームは全係数では重み一係数空間と $\mathbb A^M$ の積、
固定一次係数では $\mathbb A^M$ となる。
定理\ref{mtd:thm:normalization}の最小多項式と有限双有理射の証明を適用すれば、
既約性・被約性と正規化の主張も従う。
\end{proof}

\begin{remark}[特殊化と特異点対応の範囲]
\label{mtd:rem:scope}
全横断係数または固定 $\gamma_1$ の判別因子が既約被約であることは、
任意に固定した $(\gamma_1,\gamma_2)$ の一変数多項式が単純根だけを持つことを意味しない。
特殊 $\gamma_2$ では臨界値の衝突と非被約な臨界点を許し、
系\ref{mci:cor:finite-discriminant}の局所長を含む norm の重複度が残る。
このとき正規化射の一点上に複数の臨界点が対応する場合がある。
判別超曲面の滑らかな点は、可換標的が唯一の通常二重点を持つ係数に対応する。

ここでの通常二重点は構成された可換標的の局所型であり、
非可換中心の局所環や特異係数の全 coherent 圏の比較を結論しない。
全 coherent Ku 比較は従来どおり判別式が非零の開集合上の結論である。
本小節は固定した滑らかな境界からの逐次標的を扱い、
$W_{\mathrm{it}}$ と $W_{\mathrm{pol}}$ の同一視を用いない。
\end{remark}

% !TEX root = ../../mixed-sklyanin-cubics.tex
% Current result: 2026-11-22 / v60.
\subsection{厳密な中心作用から作る実中央 pencil と blow-up}
\label{mbl:subsec:central-blowup}

本小節では、中心方向ごとの比較核を貼り合わせる前に、source だけから
射影空間上の実圏を構成する。
内部 shift で閉じた全線束圏を用い、残余圏が中心次数付き環の全作用を
保つとは仮定しない。
構成の幾何的入力は、Artin stack に対する Khan の blow-up 分解
\cite[Theorem~C, Proposition~4.2 and Sections~4.3--4.4]{CMOKhanBlowups}である。
その分解を本稿の圏へ移す際は、実随伴と projection 三角を tensor する。
\cite[Section~1.5]{CMOJiangLeungBlowup}の幾何的 admissible 部分圏の仮定を、
本非可換代数について新たに仮定するものではない。

$k\ge2$、$N=3k$、$m\ge2$ とし、
定理\ref{mcp:thm:normal-form}の一つの複素係数を固定する。
代数を $A=A_m$、中心変数を $\mathbf u=(u_1,\ldots,u_m)$、
$S_m=\C[\mathbf u]$、$D=A/(\mathbf u)$ と書く。
$A$ は $S_m$ 上自由であり、中心列は任意順序で正則である。
$\mathcal L_A$ と $\mathcal L_D$ は、それぞれ全ての
$\OO_A(r)=\pi A(r)$、$\OO_D(r)=\pi D(r)$ の生成する冪等完備 dg 圏とする。
命題\ref{mci:prop:ambient-line-model}と固定境界の模型により、これらは有限生成かつ
$\C$-proper である。

\subsubsection{次数付き作用と中央商の実底変更}

\begin{lemma}[全線束圏の厳密な中心作用]
\label{mbl:lem:strict-action}
$\mathscr X_m=[\mathbb A^m/\mathbb G_m]$ とし、
$\mathbb G_m$ は全ての中心座標に重み一で作用する。
$\mathcal L_A$ は実 $\operatorname{Perf}(\mathscr X_m)$-線型圏である。
文字に対応する線束の作用は内部 shift であり、座標 $u_i$ の作用は
指定された中心乗算そのものである。
\end{lemma}

\begin{proof}
$\operatorname{Perf}(\mathscr X_m)$ は有限生成次数付き自由 $S_m$-加群
$S_m(r)$ の有限複体と直和因子から生成される。
このような複体 $P$ と次数付き右 $A$-加群 $M$ に対し、
$M\otimes_{S_m}^{\mathbf L}P$ は中心性によって再び右 $A$-加群となる。
$P$ を有限自由模型で表せば、これは $M$ の有限個の内部 shift と
中心乗算から作る複体である。
従って次数付き torsion 複体を保存し、qgr の導来商へ降りる。
さらに $\mathcal L_A$ を保存する。
結合子、単位、異なる $u_i$ の交換は、元の次数付き加群の tensor にある
厳密な作用から来る。
個別の自然変換の homotopy class から中心作用を作り直していない。
\end{proof}

以下の底変更では、有限 Koszul 複体を加群の解消として使う。
それが associative な相対 cofibrant 代替であるという主張は用いない。

\begin{lemma}[中央線型商の線束圏との同定]
\label{mbl:lem:central-base-change}
$R=\C[z_j\mid j\ne i]$ とし、
\[
 A_i^{\mathrm{ch}}
   =A\otimes R/(u_j-z_ju_i\mid j\ne i),\qquad
 S_i^{\mathrm{ch}}=R[u_i]
\]
と置く。$S_m\to S_i^{\mathrm{ch}}$ は $u_j\mapsto z_ju_i$ を送る。
このとき実底変更関手は同値
\begin{equation}\label{mbl:eq:chart-base-change}
 \mathcal L_A\otimes_{\operatorname{Perf}(\mathscr X_m)}
 \operatorname{Perf}([\operatorname{Spec}S_i^{\mathrm{ch}}/\mathbb G_m])
       \simeq\mathcal L_{A_i^{\mathrm{ch}}}
\end{equation}
を与える。零点についても
\begin{equation}\label{mbl:eq:center-base-change}
 \mathcal L_A\otimes_{\operatorname{Perf}(\mathscr X_m)}
             \operatorname{Perf}(B\mathbb G_m)
       \simeq\mathcal L_D
\end{equation}
である。
同じ主張は、複素数係数の独立な中央一次式で定めた任意の線型部分空間について成立する。
これらの同値は実制限関手と、有限 Koszul 解消による inflation を保持する。
\end{lemma}

\begin{proof}
まず $A$ の $S_m$-自由性から、$u_j-z_ju_i$ は $A\otimes R$ 上の正則列であり、
商は $R[u_i]$ 上自由である。
任意の複素数係数線型部分空間についても、中心座標を線型変更すれば同じである。
中心商の導来制限は線束を線束へ送り、inflation は各線束を有限 Koszul 複体へ送る。
これらは次数付き torsion を保存するので、実随伴が線束圏へ降りる。

底変更側の Hom も同じ複体で計算できる。
これは Ind 完備化における自由加群と忘却の実随伴から従う。
すなわち、基底環の商 $S'$ を代数対象として作用させた圏で、自由対象同士の Hom は
\[
 \operatorname{RHom}(F M,F M')
       =\operatorname{RHom}(M,M'\otimes^{\mathbf L}_{S_m}S')
\]
である。次数付き環上の tensor と文字の shift を保ってこの式を読む。
例えば $m=2$、$u_1=u,u_2=v$、$t=z_2$ の chart では
\begin{equation}\label{mbl:eq:one-koszul-hom}
 \begin{split}
 &\operatorname{RHom}(F\OO_A(a),F\OO_A(b))\\
 &\quad\simeq
 \operatorname{Cone}\bigl(
   \operatorname{RHom}_{\mathcal L_A}(\OO_A(a),\OO_A(b-1))[t]
    \xrightarrow{\ v-tu\ }
   \operatorname{RHom}_{\mathcal L_A}(\OO_A(a),\OO_A(b))[t]\bigr).
 \end{split}
\end{equation}
一般の chart では $m-1$ 個、零点では $m$ 個の中心射について
この Koszul totalization を取る。
通常の多項式拡大の部分では、線束が compact なので Hom が係数拡大と交換する。
既に与えた厳密な可換代数作用の自由--忘却随伴を使っており、
加群の有限解消から高次作用の有限表示を推論していない。

一方、実 qgr 制限と inflation の随伴によって
\[
 \operatorname{RHom}_{\mathcal L_{A_i^{\mathrm{ch}}}}
       (\OO(a),\OO(b))
  =\operatorname{RHom}_{\mathcal L_{A\otimes R}}
       (\OO(a),i_*\OO(b))
\]
であり、右辺は同じ Koszul 複体である。
中心作用を保つ実制限から式\eqref{mbl:eq:chart-base-change}の関手が定まり、
上の計算はその Hom 上の射が quasi-isomorphism であることを示す。
基底 stack の perfect 圏は文字付き自由加群で生成され、
右辺は全線束で生成されるので、この関手は Morita 同値である。
零点と一般の線型部分空間も同じ証明による。
比較は実関手上で行ったため、積、随伴射、高次 coherence も保持される。
\end{proof}

\subsubsection{射影空間上の実圏と各 fiber}

\[
 \mathscr Y_m=[\operatorname{Bl}_0\mathbb A^m/\mathbb G_m],\qquad
 b:\mathscr Y_m\longrightarrow\mathscr X_m,\qquad
 p:\mathscr Y_m\longrightarrow\mathbb P^{m-1}
\]
と置く。後者の射は、例外因子を含む射影方向を取る射である。

\begin{theorem}[全中心 pencil の実線束圏]
\label{mbl:thm:pencil-category}
\begin{equation}\label{mbl:eq:pencil-category}
 \mathcal P_A
   =\mathcal L_A\otimes_{\operatorname{Perf}(\mathscr X_m)}
                       \operatorname{Perf}(\mathscr Y_m)
\end{equation}
は実 $\mathbb P^{m-1}$-線型圏である。
標準 chart $U_i=\{\ell_i\ne0\}$ では
\[
                  \mathcal P_A|_{U_i}\simeq\mathcal L_{A_i^{\mathrm{ch}}}.
\]
従って任意の点 $\ell=[\ell_1:\cdots:\ell_m]$ における導来 fiber は
\begin{equation}\label{mbl:eq:ray-fiber}
  (\mathcal P_A)_\ell
   \simeq\mathcal L_{A_\ell},\qquad
  A_\ell=A/(\ell_i u_j-\ell_j u_i\mid j\ne i),
       \quad\ell_i\ne0
\end{equation}
である。右辺は一中心商の通常の線束圏であり、追加の導来厚みを持たない。
$\mathcal P_A$ は $\mathbb P^{m-1}$ 上 proper である。
\end{theorem}

\begin{proof}
$p$ に沿う引戻しにより、式\eqref{mbl:eq:pencil-category}に
$\operatorname{Perf}(\mathbb P^{m-1})$ の作用が入る。
blow-up の $i$ 番目の chart は
\[
 [\operatorname{Spec}\C[u_i,z_j\mid j\ne i]/\mathbb G_m],
 \qquad \deg(u_i)=1,\quad\deg(z_j)=0,
\]
であり、$b$ は $u_j=z_ju_i$ で与えられる。
perfect 圏のこの chart への底変更は、次数付き加群の開局所化で計算できる。
従って tensor の結合と補題\ref{mbl:lem:central-base-change}により chart の同定を得る。
重なりでは、例えば二中心の場合
\[
 A[t,t^{-1}]/(v-tu)
       \simeq A[s,s^{-1}]/(u-sv),\qquad s=t^{-1}
\]
という元の $A$ 上の次数付き同型を用いる。
一般の重なりも同じ座標比の同型であり、三重交叉上の coherence は
元の incidence chart のものそのままである。
個々の ray の Morita 核を選んで貼り合わせているのではない。

$A_i^{\mathrm{ch}}$ は $\C[\mathbf z,u_i]$ 上自由なので、
$\mathbf z=\boldsymbol\lambda$ への底変更後も表示した通常商が得られる。
Hom の Koszul 模型もこの底変更と交換する。
再び線束生成を使うと式\eqref{mbl:eq:ray-fiber}を得る。

各 chart の代数は、中心形式を
$\mathbf u=u_i(\ell_1(\mathbf z),\ldots,\ell_m(\mathbf z))$ とした一中心族である。
命題\ref{emc:prop:line-model}の全係数模型と任意底変更を用いると、
その線束間 Hom は $\C[\mathbf z]$ 上 perfect となる。
これは直接にも、相対線束の有限 Koszul・局所 cohomology 模型から従う。
$\mathcal P_A$ の任意の二対象の $\mathbb P^{m-1}$-値 Hom は、
chart への底変更でこれらの Hom に一致する。
mapping object のこの底変更は自由対象上の計算と有限 cone・直和因子から従い、
\cite[Lemmas~2.7 and 2.10]{CRSPerry}の形式に一致する。
perfect 性は底の Zariski 局所で判定できるので、相対 proper 性を得る。
\end{proof}

この chart は通常の graded qgr としても具体的に書ける。
射影方向の斉次座標を $\ell_i$ とすると、例えば二中心では、各次数の加群は
\[
 \mathscr A_n
   =\operatorname{Coker}\bigl(
    A_{n-1}\otimes\OO_{\mathbb P^1}(-1)
       \xrightarrow{\ \ell_1v-\ell_2u\ }
    A_n\otimes\OO_{\mathbb P^1}\bigr).
\]
各点で中央一次式が正則なので、これは vector bundle である。
中心性により次数付き代数の積が入り、二つの chart で上の通常商を与える。
定理の tensor 定義は、こうした通常 chart と実随伴の情報を同時に保持する。

\subsubsection{Orlov 分解、proper 性と絶対 smoothness}

例外因子を $j:E_m\hookrightarrow\mathscr Y_m$、その射を
$q:E_m\to B\mathbb G_m$ と書く。
$E_m\simeq\mathbb P^{m-1}\times B\mathbb G_m$ だが、
線束 $\OO_{E_m}(-1)=\OO_{\mathscr Y_m}(E_m)|_{E_m}$ には元の文字を保持する。
内部次数を忘れた $\mathbb P^{m-1}$ の線束へ置き換えない。

\begin{proposition}[中心 blow-up の実半直交分解]
\label{mbl:prop:orlov}
実完全忠実関手 $b^*:\mathcal L_A\to\mathcal P_A$ と
$\Phi_r:\mathcal L_D\to\mathcal P_A$ $(1\le r\le m-1)$ があり、
\begin{equation}\label{mbl:eq:orlov}
 \mathcal P_A
   =\left\langle
      \Phi_{m-1}\mathcal L_D,\ldots,\Phi_1\mathcal L_D,
      b^*\mathcal L_A
     \right\rangle
\end{equation}
となる。各関手は compact 対象を保つ左右随伴を持つ。
幾何側の $\Phi_r$ は $j_*(q^*(-)\otimes\OO_{E_m}(-r))$ である。
特に $\mathcal P_A$ は $\C$ 上 proper かつ有限生成であり、
$\mathcal L_A,\mathcal L_D$ がともに smooth なら $\mathcal P_A$ も smooth である。

二中心の場合、$i^*:\mathcal L_A\to\mathcal L_D$ を中心列全体への制限とすると、
gluing は実同型
\begin{equation}\label{mbl:eq:gluing}
 \begin{split}
 \operatorname{RHom}_{\mathcal P_A}(b^*M,\Phi_1 N)&=0,\\
 \operatorname{RHom}_{\mathcal P_A}(\Phi_1 N,b^*M)
      &\simeq\operatorname{RHom}_{\mathcal L_D}(N,i^*M)[-1]
 \end{split}
\end{equation}
で与えられる。
\end{proposition}

\begin{proof}
$B\mathbb G_m\to\mathscr X_m$ は正則な余次元 $m$ の閉埋込みである。
Khan の定理\cite[Theorem~C]{CMOKhanBlowups}は、これらの Artin stack の
perfect 圏に式\eqref{mbl:eq:orlov}の形の分解を与える。
同文献の表示順序は本稿の半直交分解の規約と逆なので、ここでは
$\operatorname{RHom}(b^*(-),j_*(q^*(-)\otimes\OO(-r)))=0$
となる順序を採用した。

この分解は $\operatorname{Perf}(\mathscr X_m)$-線型である。
引戻し、proper pushforward、$j^!$、および必要な線束 twist は perfect 対象を保つ。
右随伴は、例えば
\[
 b_*,\qquad q_*(j^!(-)\otimes\OO_{E_m}(r))
\]
であり、左随伴も relative dualizing line bundle による双対から得られる。
これらは実随伴の単位・余単位と projection 三角を持つ。
それを $\mathcal L_A$ と相対 tensor すれば、単位の同型、直交性、
有限 projection 三角が保たれる。
これは\cite[Lemma~3.15]{CRSPerry}の証明と同じ形式論であるが、
ここで必要なのは底の perfect 圏が rigid symmetric monoidal であることと
上の実線型随伴であり、底をスキームとする仮定ではない。
従って scheme 版の文言を stack へ無条件に読み替えていない。
補題\ref{mbl:lem:central-base-change}によって零点成分が $\mathcal L_D$ となり、
式\eqref{mbl:eq:orlov}と実随伴を得る。

gluing の有限性を確認する。
Cartier 双対\cite[Proposition~4.2]{CMOKhanBlowups}は
\[
 j^!b^*M=q^*i^*M\otimes\OO_{E_m}(-1)[-1]
\]
を与える。$m=2$ では $Rq_*\OO_{E_2}=\OO$、
$Rq_*\OO_{E_2}(-1)=0$ なので式\eqref{mbl:eq:gluing}を得る。
特に gluing は $\C$ 上 perfect である。
一般 $m$ でも、$\Phi_r$ の compact 右随伴を用いると全 cross Hom は
\[
 \operatorname{RHom}_{\mathcal L_D}(N,\Phi_r^R\Phi_sN')
 \quad\text{または}\quad
 \operatorname{RHom}_{\mathcal L_D}(N,\Phi_r^Rb^*M)
\]
となり、$\mathcal L_D$ の proper 性から perfect である。
$\mathcal L_A,\mathcal L_D$ の有限生成対象の像を合わせれば $\mathcal P_A$ を生成するので、
絶対 proper 性と有限生成性が従う。
相対 proper 性と $\mathbb P^{m-1}$ の proper 性からも同じ結論が得られる。

最後に二つの成分が smooth であるとする。
その有限生成代数を用いると、gluing 双加群は $\C$ 上 perfect であり、
両側の smooth 対角の有限解消を tensor することで双加群としても perfect になる。
有限三角 dg 代数の対角は各成分の対角とこの gluing から有限 cone で作られる。
従って $\mathcal P_A$ は smooth である。
この最後の段階は単に数値不変量が加法的であるという議論ではない。
\end{proof}

\begin{corollary}[滑らかな source の実 pencil]
\label{mbl:cor:smooth-source}
本稿の固定楕円境界では $\mathcal L_D$ は smooth proper である。
従って $\mathcal L_A$ が smooth なら $\mathcal P_A$ は絶対的に smooth proper である。
特に $A$ の残余線束圏が smooth ならこの結論が成立する。
この段階では、全ての fiber が smooth、あるいは一般 fiber が smooth であるとはまだ結論しない。
\end{corollary}

\begin{proof}
補題\ref{ebk:lem:boundary}により $D$ の残余線束圏は滑らかな可換三次の Ku と同値である。
有限例外線束を戻す gluing は proper であり、前証明と同じ有限対角の議論から
$\mathcal L_D$ が smooth となる。
同じ理由で、$A$ の残余線束圏が smooth なら $\mathcal L_A$ も smooth である。
命題\ref{mbl:prop:orlov}を適用する。
\end{proof}

\subsubsection{中央線型切断との適合と実 chart 比較}

\begin{lemma}[中心 hyperplane への制限]
\label{mbl:lem:linear-section}
$m\ge3$ とし、中心線型形式 $\lambda$ に対応する超平面
$H=\mathbb P(\ker\lambda)\subset\mathbb P^{m-1}$ を取る。
$A'=A/(\lambda)$ とすると、実同値
\begin{equation}\label{mbl:eq:linear-section}
        \mathcal P_A\otimes_{\operatorname{Perf}(\mathbb P^{m-1})}
                       \operatorname{Perf}(H)
              \simeq\mathcal P_{A'}
\end{equation}
がある。
右辺は $m-1$ 中心の同じ構成であり、
$m-2$ 個の $\mathcal L_D$ と $\mathcal L_{A'}$ による実分解を持つ。
$m=2$ では左辺は式\eqref{mbl:eq:ray-fiber}の一中心商の線束圏である。
\end{lemma}

\begin{proof}
中心座標を変えて $\lambda=u_m$ としてよい。
$\mathscr X_H=[\ker\lambda/\mathbb G_m]$ と置くと、
\[
 \mathscr Y_m|_H=[\operatorname{Tot}\OO_H(-1)/\mathbb G_m]
               =\operatorname{Bl}_{B\mathbb G_m}\mathscr X_H.
\]
$\mathscr Y_m\to\mathbb P^{m-1}$ は線束の全空間の商として平坦であり、
この base change は通常のものと一致する。
その $\mathscr X_m$ への射は $\mathscr X_H$ を通る。
従って tensor の結合と補題\ref{mbl:lem:central-base-change}から
\[
 \begin{split}
 \mathcal P_A|_H
  &\simeq
  \bigl(\mathcal L_A\otimes_{\operatorname{Perf}(\mathscr X_m)}
                      \operatorname{Perf}(\mathscr X_H)\bigr)
          \otimes_{\operatorname{Perf}(\mathscr X_H)}
                      \operatorname{Perf}(\mathscr Y_m|_H)\\
  &\simeq\mathcal P_{A'}
 \end{split}
\]
となる。ここで $\lambda$ は正則なので、$A'$ は導来商でなく通常の商である。
chart では $z_m=0$ とするだけであり、実 restriction と随伴もこの同定を保つ。
\end{proof}

\begin{corollary}[一中心普遍核の各 chart への実適用]
\label{mbl:cor:chart-kernel}
$\mathcal P_A$ には $\mathbb P^{m-1}$-線型な実分解
\[
 \mathcal P_A
  =\left\langle\mathcal K_{\mathrm{pen}},
     \OO\otimes\operatorname{Perf}(\mathbb P^{m-1}),\ldots,
     \OO(N-3)\otimes\operatorname{Perf}(\mathbb P^{m-1})\right\rangle
\]
がある。$\mathcal K_{\mathrm{pen}}|_{U_i}$ は $A_i^{\mathrm{ch}}$ の残余線束圏である。

一中心の定理\ref{euk:thm:universal-kernel}で係数多項式 $Q,L,C$ と実核を固定する。
中心形式を代入した
\[
 \begin{split}
 W_{\mathrm{pol}}(\mathbf x,\mathbf u)
   ={}&G(\mathbf x)+Q(\mathbf t(\mathbf u);\mathbf x)\\
     &+L(\mathbf t(\mathbf u),\alpha(\mathbf u);\mathbf x)
      +C(\mathbf t(\mathbf u),\alpha(\mathbf u),c(\mathbf u))
 \end{split}
\]
を用い、$\ell_i(\mathbf z)=1$ とする chart の方向ベクトルを
$\boldsymbol\ell^{(i)}(\mathbf z)$ と書く。
このとき $\mathcal K_{\mathrm{pen}}|_{U_i}$ は、全 chart 係数環 $\C[\mathbf z]$ 上で
\[
       V\bigl(W_{\mathrm{pol}}(\mathbf x,s\boldsymbol\ell^{(i)}(\mathbf z))\bigr)
            \subset\mathbb P^N\times U_i
\]
の残余線束圏と実 Morita 同値である。
この同値は任意の導来底変更と両立する。
\end{corollary}

\begin{proof}
$\OO(r)=b^*\OO_A(r)$ と書く。
chart の相対線束 Hom により、$\OO,\ldots,\OO(N-3)$ は
$\mathbb P^{m-1}$ 上の相対例外列である。
相対 proper 性から、その左右随伴はそれぞれ相対 Hom とその双対で与えられ、
perfect 対象を保つ。
有限 mutation で残余 $\mathcal K_{\mathrm{pen}}$ を定めれば、この構成は各 chart と
底変更に両立する。
元の $\mathcal K_A$ が内部 shift に閉じているとは仮定していない。

$A_i^{\mathrm{ch}}$ は、一中心の全係数環から
\[
 \mathbf t\longmapsto\mathbf t(\boldsymbol\ell^{(i)}),\qquad
 \alpha\longmapsto\alpha(\boldsymbol\ell^{(i)}),\qquad
 c\longmapsto c(\boldsymbol\ell^{(i)})
\]
という多項式係数変更で得られる。
中心持上げも補題\ref{mcp:lem:polarization}によってこの代入を保つ。
重み $(1,2,3)$ から、可換側の三次は正確に
$W_{\mathrm{pol}}(\mathbf x,s\boldsymbol\ell^{(i)})$ となる。
従って一中心の実普遍核をこの係数環へ導来底変更すればよい。
双方の次数別自由性により、ここでの代数は通常の係数変更と一致する。
\end{proof}

\begin{remark}[構成した範囲]
\label{mbl:rem:scope}
本小節で得たのは、厳密な source 中心作用から作る実射影 pencil、その各通常 chart、
実 blow-up 分解、および source が smooth な場合の全 pencil 圏の絶対 smoothness である。
さらに各 chart の残余比較は、既存の一中心普遍核の実底変更として存在する。
異なる chart の非可換・可換比較核が重なりで同時に降下するとは主張しない。
個別の比較核に二つ以上の中心作用を coherent に同時保存させる問題とも区別する。
一般 fiber の smoothness と良い中央 flag の存在は、
次の定理\ref{mcb:thm:smooth-locus}と定理\ref{mcb:thm:central-flags}で
この実構成に続いて証明する。
\end{remark}

% !TEX root = ../../mixed-sklyanin-cubics.tex
% Current result: 2026-11-22 / v60.
\subsection{偏極三次による全 smooth locus と良い中央旗}
\label{mcb:subsec:central-bertini}

前小節の実 blow-up と、一中心普遍核の各 chart への適用を組み合わせると、
中央変数の個数が任意の場合も smoothness を判定できる。
用いる局所性は、$\C$ 上の\emph{絶対対角}の perfect 性である。
射影方向上の相対 smoothness や、異なる chart の比較核の降下は仮定しない。

\subsubsection{絶対対角の Zariski 局所性と相対 gluing}

\begin{lemma}[絶対 smoothness の局所判定]
\label{mcb:lem:smooth-local}
$S$ を滑らかな有限型分離 $\C$-スキームとし、
$\mathcal C$ を実 $\operatorname{Perf}(S)$-線型の冪等完備 dg 圏とする。
$S$ の有限 affine 開被覆 $S=\bigcup_i U_i$ に対し
\[
 \mathcal C_{U_i}
   =\mathcal C\otimes_{\operatorname{Perf}(S)}\operatorname{Perf}(U_i)
\]
と置く。このとき、$\mathcal C$ が $\C$ 上 smooth であることと、
全ての $\mathcal C_{U_i}$ が $\C$ 上 smooth であることは同値である。
ここでは $S$ 上の相対対角の perfect 性を結論しない。
また、両辺の圏が $\C$ 上 proper であることは必要ない。
\end{lemma}

\begin{proof}
有限開被覆による compact 性の判定を先に記す。
$T$ を有限型分離 $\C$-スキーム、$\mathcal Q$ を
$\operatorname{QCoh}(T)$-線型の presentable stable dg 圏とする。
affine 開埋込み $j:V\hookrightarrow T$ に対し、
$j_*\OO_V$ は $\operatorname{QCoh}(T)$ の idempotent algebra である。
その作用による局所化を $j^*:\mathcal Q\to\mathcal Q_V$ と書くと、
忘却は完全忠実な右随伴 $j_*$ であり、
\[
             j_*j^*N=(j_*\OO_V)\otimes_{\OO_T}N
\]
となる。この右随伴は filtered colimit を保つので、$j^*$ は compact 対象を保つ。

有限 affine 開被覆 $\{V_a\}_{a=1}^r$ を取る。
分離性から全ての有限交叉 $V_I=\bigcap_{a\in I}V_a$ は affine であり、
$j_I:V_I\hookrightarrow T$ も affine な開埋込みである。
構造層の有界 augmented \v{C}ech 複体を $N$ に作用させると
\begin{equation}\label{mcb:eq:cech}
 N\simeq\operatorname{Tot}\left[
   \bigoplus_{|I|=1}j_{I,*}j_I^*N\longrightarrow
   \bigoplus_{|I|=2}j_{I,*}j_I^*N\longrightarrow\cdots
   \longrightarrow j_{\{1,\ldots,r\},*}j_{\{1,\ldots,r\}}^*N
                         \right].
\end{equation}
空な交叉の項は零とする。この totalization は有限である。
各 $j_a^*M$ が compact なら交叉への制限も compact であり、随伴により
\[
 \operatorname{Map}_{\mathcal Q}(M,j_{I,*}j_I^*N)
   =\operatorname{Map}_{\mathcal Q_{V_I}}(j_I^*M,j_I^*N)
\]
となる。filtered colimit と有限 totalization が交換するので、
式\eqref{mcb:eq:cech}は $M$ 自身が compact であることを示す。
これは線型圏の descent における有限 mapping pullback の議論
\cite[Lemma~6.4 and proof of Theorem~6.11, pp.~1222--1223]{CMOAntieauGepnerDescent}
を、ここで必要な有限開被覆について直接書いたものである。

これを
\[
 T=S\times_\C S,\qquad
 \mathcal Q=\operatorname{Ind}
                (\mathcal C^{\mathrm{op}}\otimes_\C\mathcal C)
\]
に適用する。二つの既存の $\operatorname{Perf}(S)$ 作用は
$\operatorname{Perf}(S\times S)$ 作用を与え、Ind 完備化後に
$\operatorname{QCoh}(S\times S)$ へ延長される。
$M_{\mathcal C}\in\mathcal Q$ を絶対対角双加群とする。
恒等関手 $\operatorname{id}_{\mathcal C}$ は、元の実作用に関して
canonically $S$-線型である。
従ってその graph は相対 envelope 上の対角からの restriction として現れ、
二つの底作用は $\Delta_S:S\to S\times S$ に沿って一致する。
特に $M_{\mathcal C}$ の $S\times S\setminus\Delta_S$ への制限は零である。
affine 底では、これは対角上の左右のスカラー作用が同一であり、
$a\otimes1-1\otimes a$ が零作用するという通常の双加群の事実である。
その全 coherence は、既に与えた $\operatorname{Perf}(S)$ 作用と恒等関手から来る。
外部から与えた Morita 核に別々の中心作用を新たに整合させる問題は生じない。

一方、$U_i\times U_i$ への制限は
$\mathcal C_{U_i}$ の絶対対角である。
実際、開制限は Morita 局所化であり、左右から誘導した対角は局所化後の対角となる。
$\{U_i\times U_i\}_i$ と $S\times S\setminus\Delta_S$ は $S\times S$ を覆う。
後者を有限 affine 開で覆い、全体を有限 affine 被覆にすれば、
前段の compact 性の局所判定が使える。
従って全 $\mathcal C_{U_i}$ が smooth なら $M_{\mathcal C}$ は compact となる。
逆向きは、開局所化が compact 対象と対角を保つことから従う。
\end{proof}

\begin{lemma}[相対例外列を戻すときの絶対 smoothness]
\label{mcb:lem:relative-gluing}
$R=\C[t_1,\ldots,t_r]$ とし、$\mathcal C$ を $R$-proper な有限生成 dg 圏とする。
実 admissible 半直交分解
\[
 \mathcal C=\langle\mathcal K,
       E_1\otimes\operatorname{Perf}(R),\ldots,
       E_a\otimes\operatorname{Perf}(R)\rangle
\]
を持つと仮定する。ここで $E_j$ は相対例外対象であり、
全ての随伴は compact 対象を保つものとする。
このとき $\mathcal C$ が $\C$ 上 smooth であることと、
$\mathcal K$ が $\C$ 上 smooth であることは同値である。
\end{lemma}

\begin{proof}
admissible 部分圏の絶対対角は、周囲の対角を左右随伴によって制限して得られる。
随伴が compact 対象を保つため、$\mathcal C$ の smoothness は $\mathcal K$ へ移る。

逆に $\mathcal K$ が smooth とする。
その有限生成代数を $E$ とすると、最初の例外成分との cross Hom は
$E$ と $R$ の双加群 $M$ であり、$R$-proper 性から underlying $R$-複体が perfect である。
$E$ の絶対対角の有限 $E^{\mathrm{op}}\otimes_\C E$-perfect 解消を
片側で $M$ と tensor する。
自由な対角項は $E\otimes_\C M$ となり、これは
$E\otimes_\C R$ 上 perfect である。
従って $M$ は、向きを適切に取った
$E^{\mathrm{op}}\otimes_\C R$ 上の perfect 双加群となる。
$R$ 自身も $\C$ 上 smooth なので、三角 dg 代数の対角解消
\cite[Proposition~3.11]{CMOLuntsResolution}により、
最初の例外成分を戻した圏は絶対 smooth である。
この圏は引き続き $R$-proper であるから、同じ議論を有限回反復できる。
smooth な成分だけから gluing の perfect 性を省略しているのではなく、
その不足を $R$-proper 性で補っている。
\end{proof}

\subsubsection{全多中心族の smoothness の必要十分条件}

定理\ref{mcp:thm:normal-form}の固定複素係数に戻り、$m\ge1$ とする。
一中心の比較で一度固定した係数多項式 $Q,L,C$ に中心形式を代入し、
\begin{equation}\label{mcb:eq:polarized}
 \begin{split}
 W_{\mathrm{pol}}(\mathbf x,\mathbf u)
  ={}&G(\mathbf x)+Q(\mathbf t(\mathbf u))(\mathbf x)\\
    &+L(\mathbf t(\mathbf u),\alpha(\mathbf u))(\mathbf x)
      +C(\mathbf t(\mathbf u),\alpha(\mathbf u),c(\mathbf u)),\\
 X_{\mathrm{pol}}={}&V(W_{\mathrm{pol}})\subset\mathbb P^{N+m-1},\qquad
 Z=V(G)\subset\{\mathbf u=0\}\simeq\mathbb P^{N-1}
 \end{split}
\end{equation}
と定める。$m=1$ は定理\ref{euk:thm:universal-kernel}、
$m=2$ は式\eqref{mcpn:eq:polarized-candidate}の三次である。
$Z$ は滑らかであり、$Z$ の各点では $\partial G/\partial x_j$ の少なくとも一つが非零である。
従って $X_{\mathrm{pol}}$ は係数によらず $Z$ の近傍で滑らかである。

\begin{theorem}[全多中心族の smooth locus]
\label{mcb:thm:smooth-locus}
$k\ge2$、$N=3k$ とし、任意の固定した滑らかな標準楕円 Sklyanin 因子、
非零重み、および定理\ref{mcp:thm:normal-form}の任意の複素係数を取る。
$\mathcal K_A$ を $A=A_m$ の残余線束圏とすると
\begin{equation}\label{mcb:eq:smooth-locus}
       \mathcal K_A\text{ が $\C$ 上 smooth}
          \quad\Longleftrightarrow\quad
       X_{\mathrm{pol}}\text{ が滑らか}
\end{equation}
である。二捩れ因子と非二捩れ因子の混在、中心変数の個数、
係数の一般性について追加条件は要らない。
また、代数が有限型中央部分環上有限であるという条件も要らない。
\end{theorem}

\begin{proof}
$m=1$ は既存の一中心比較である。以下 $m\ge2$ とする。
全線束圏 $\mathcal L_A$ は proper であり、
$\mathcal K_A$ と有限例外線束による実分解を持つ。
従って補題\ref{mcb:lem:relative-gluing}の $R=\C$ の場合により、
$\mathcal K_A$ と $\mathcal L_A$ の絶対 smoothness は同値である。
固定境界の $\mathcal L_D$ は smooth proper なので、
命題\ref{mbl:prop:orlov}と $b^*\mathcal L_A$ の admissibility から
\[
             \mathcal L_A\text{ smooth}
                \quad\Longleftrightarrow\quad
             \mathcal P_A\text{ smooth}
\]
を得る。

可換側の blow-up を $\widetilde X=\operatorname{Bl}_Z X_{\mathrm{pol}}$ とする。
その $\mathbb P^{m-1}$ 上の標準 chart は
\[
 X_i=V\bigl(W_{\mathrm{pol}}
       (\mathbf x,s\boldsymbol\ell^{(i)}(\mathbf t))\bigr)
                    \subset\mathbb P^N\times U_i
\]
である。ここで $\ell_i^{(i)}=1$ とする。
この chart 同定は特異な係数でも有効である。
実際、$G\ne0$ なので $(u_1,\ldots,u_m)$ は三次商上の正則列であり、
表示した chart の $s=0$ への制限は $G=0$ である。
従って $s$ は非零因子であり、incidence chart には例外因子にのみ支持された
余分な成分はない。これは rational map $[\mathbf x:\mathbf u]\mapsto[\mathbf u]$
の graph の閉包、すなわちこの ideal の blow-up を与える。

系\ref{mbl:cor:chart-kernel}は、各 $R_i=\C[\mathbf t]$ 上で
$\mathcal P_A|_{U_i}$ の残余と $X_i$ の残余線束圏との実 Morita 同値を与える。
両側の相対線束模型は $R_i$-proper であり、相対例外列の左右随伴は compact を保つ。
また $\operatorname{Perf}(X_i)$ は全相対線束 twist で生成される。
これは射影スキームの ample 線束による生成であり、底 $U_i$ が affine なので
$\operatorname{Perf}(R_i)$ の係数も有限 cone と直和因子で含まれる。
従って全 coherent 圏が先に飽和していると仮定せず、
補題\ref{mcb:lem:relative-gluing}を両側に適用して
\begin{equation}\label{mcb:eq:chart-smooth}
 \mathcal P_A|_{U_i}\text{ が $\C$ 上 smooth}
   \quad\Longleftrightarrow\quad
 \operatorname{Perf}(X_i)\text{ が $\C$ 上 smooth}
   \quad\Longleftrightarrow\quad X_i\text{ が滑らか}
\end{equation}
を得る。最後の同値は、分離有限型スキームに対する
\cite[Proposition~3.13]{CMOLuntsResolution}である。
$X_i$ は $\C$ 上 proper とは限らないが、同命題に proper 性は必要ない。

補題\ref{mcb:lem:smooth-local}を $S=\mathbb P^{m-1}$ と
$\mathcal C=\mathcal P_A$ に適用する。
式\eqref{mcb:eq:chart-smooth}と通常のスキームの smoothness の局所性により
\[
       \mathcal P_A\text{ smooth}
          \quad\Longleftrightarrow\quad\widetilde X\text{ smooth}
\]
となる。$X_{\mathrm{pol}}$ が滑らかなら、滑らかな中心 $Z$ の blow-up も滑らかである。
逆に $\widetilde X$ が滑らかなら、blow-up が同型である $X_{\mathrm{pol}}\setminus Z$ は滑らかである。
$Z$ の近傍は既に $\partial G/\partial x_j$ によって滑らかと分かっている。
これで式\eqref{mcb:eq:smooth-locus}の両方向が従う。
\end{proof}

\subsubsection{一般中央超平面、完全旗と全 coherent 比較}

\begin{theorem}[滑らかな対象の一般中央旗]
\label{mcb:thm:central-flags}
定理\ref{mcb:thm:smooth-locus}の $A_m$ について、
$\mathcal K_{A_m}$ が smooth であると仮定する。
中心空間の双対射影空間には非空 Zariski 開集合があり、
その任意の中央一次式 $\lambda$ に対して $A_m/(\lambda)$ の残余線束圏は smooth である。

さらに、中心空間の完全旗のパラメータ空間には非空 Zariski 開集合があり、
その旗に沿う全ての通常商
\[
 A_j=A_m/(\lambda_{j+1},\ldots,\lambda_m),\qquad 1\le j\le m,
\]
が smooth な残余線束圏を持つ。
ここで $\lambda_1,\ldots,\lambda_m$ は旗に適合した中心座標であり、
$A_0=D$ は固定された滑らかな境界である。
\end{theorem}

\begin{proof}
$m=1$ では中央一次式による商は固定境界 $D$ であり、旗に中間段はない。
その smoothness は既知である。以下 $m\ge2$ とする。
定理\ref{mcb:thm:smooth-locus}により $X=X_{\mathrm{pol}}$ は滑らかである。
中央一次式からなる線型系の基点は $Z$ である。
以下で、この基点付き線型系に必要な古典的 Bertini の議論を直接確認する。

$V=\C^m$ とし、普遍中央超平面を
\[
 \mathscr H=\{(x,[\lambda])\in X\times\mathbb P(V^*)
                               \mid\lambda(\mathbf u(x))=0\}
\]
と置く。$X\setminus Z$ 上では係数 $\lambda$ 方向の微分により $\mathscr H$ は滑らかである。
$Z$ 上では $Z\hookrightarrow X$ は余次元 $m$ の滑らかな正則埋込みであり、
$du_1,\ldots,du_m$ は conormal 方向で独立である。
従って非零の $\lambda$ について $\sum_i\lambda_i\,du_i$ は非零となり、
そこでも $\mathscr H$ は滑らかである。

射影 $\mathscr H\to\mathbb P(V^*)$ は proper である。
その generic fiber の局所環は滑らかな全空間の局所環の局所化なので regular である。
標数零の関数体は perfect であり、この有限型 generic fiber は smooth となる。
非 smooth locus は閉集合で generic fiber と交わらない。
proper 像を除けば非空開集合上の全ファイバーが smooth となる。
これが必要な一般中央超平面の存在である。

中心座標の線型変更は定理\ref{mcp:thm:normal-form}と自由性を保つ。
また $Q,L,C$ は係数多項式への代入で定義したので、
$A_m/(\lambda)$ に対応する偏極三次は正確に $W_{\mathrm{pol}}|_{\lambda=0}$ である。
従って定理\ref{mcb:thm:smooth-locus}を $m-1$ 中心の商へ適用すれば、
その残余線束圏も smooth である。$m-1=1$ では一中心の同値を用いる。
これを帰納して良い旗を一つ得る。
各段の普遍可換切断は、非零の固定項 $G$ を持つ相対三次超曲面として
旗空間上 flat かつ projective であり、smooth fiber の条件は開である。
有限個の段についてその開集合を交差すれば、表示した非空開集合を得る。
全商が通常の商であることは、任意順序の中心正則性による。
\end{proof}

\begin{corollary}[二中心の良い ray と判別形式]
\label{mcb:cor:good-ray}
$m=2$ で $\mathcal K_A$ が smooth なら、
命題\ref{mcpn:prop:ray-discriminant}の斉次形式 $\Delta$ は恒等零ではない。
良い中央切断は正確に $\mathbb P^1\setminus V(\Delta)$ であり、
悪い方向は有限個である。
特に $\Delta\equiv0$ の係数には、smooth な残余線束圏を持つ対象は存在しない。
\end{corollary}

\begin{proof}
定理\ref{mcb:thm:central-flags}が少なくとも一つの smooth な一中心商を与える。
命題\ref{mcpn:prop:ray-discriminant}の必要十分条件を適用する。
この結論は $\Delta$ の全ての根が単純であるとは主張しない。
\end{proof}

\begin{corollary}[全 factorwise smooth 対象の可換 Ku 比較]
\label{mcb:cor:all-smooth-comparison}
定理\ref{mcb:thm:smooth-locus}の因子別多中心族で、
残余線束圏が smooth な任意の $A_m$ は、
ある滑らかな可換三次
\[
                       Y_{\mathrm{it}}\subset\mathbb P^{N+m-1}
\]
を持ち、$\dim Y_{\mathrm{it}}=N+m-2$ で、全 coherent 残余圏について
\[
             \operatorname{Ku}_{\mathrm{coh}}(A_m)
                   \simeq\operatorname{Ku}(Y_{\mathrm{it}})
\]
という実 Morita 同値がある。
良い中央旗を選べば、各段ごとの実制限随伴、単位・余単位と周期を保つ逐次比較を取れる。
比較核は対応する smooth な圏の積上 perfect である。
\end{corollary}

\begin{proof}
$m=1$ は定理\ref{euk:thm:universal-kernel}の smooth locus における全 coherent 比較である。
$m\ge2$ では
定理\ref{mcb:thm:central-flags}の旗を固定し、
系\ref{mci:cor:smooth-flag}の逐次構成を適用する。
各境界は smooth なので、必要な実比較と同期した周期が次段の入力となる。
最終段も smooth であり、同系の飽和論証が全 coherent 残余圏への同値と
perfect な graph 核を与える。
\end{proof}

\begin{remark}[判定と核の適用範囲]
\label{mcb:rem:scope}
式\eqref{mcb:eq:smooth-locus}は、偏極三次による\emph{全 smooth locus の判定}である。
逐次構成の三次式 $W_{\mathrm{it}}$ と $W_{\mathrm{pol}}$ の一致、
あるいは $\mathcal K_A$ と $\operatorname{Ku}(X_{\mathrm{pol}})$ の同値を主張するものではない。
各 chart の核は smoothness の移送に用いたが、核自体を全射影空間上へ貼り合わせてはいない。
旗に沿う比較で保存する随伴と周期は各段ごとのものであり、
以前の全ての中心作用の標識を最終核で同時に保存するという結論ではない。
逐次同値が保つのは各段ごとの実制限随伴と周期であり、
以前の段の全中心標識を一つの核で同時に保存するという意味ではない。
また、一般の smooth proper dg 圏に対する Bertini 定理を証明したのでもない。
本結果では、実中心作用から作った blow-up と、全係数上の一中心普遍核が、
絶対 smoothness を通常の有限型スキームへ移すための具体的な入力を与えている。
\end{remark}

% !TEX root = ../../mixed-sklyanin-cubics.tex
% Current result: 2026-11-22 / v59.
\subsection{二重制限の実随伴正方形と混合構造の必要性}
\label{tcr:subsec:restriction-square}

二つの中心方向の制限は、残余圏上でも随伴に適合して交換する。
しかし、二つの軸の spherical comonad をそれぞれ同定しても、
全対象の圏は決まらない。本小節では、この二つを区別する。
前半は正則列と mutation による実 Beck--Chevalley 射の計算であり、
後半は混合項の有限行列計算と、軸が同一でも全圏の smoothness が異なる例である。

$k\ge2$、$N=3k$ とし、定理\ref{mcp:thm:normal-form}の
固定した複素係数、二つの中心変数 $u,v$ を持つ代数を $A$ と書く。
\[
 A_u=A/(u),\qquad A_v=A/(v),\qquad S=A/(u,v)
\]
と置く。可換三次の場合も同じ記号を用いる。
各代数 $R$ の全線束 twist の生成圏を $\mathcal C_R$、
標準例外線束列の右直交を $\mathcal K_R$ とする。
例外線束列の長さはそれぞれ
\[
 r_A=N-1,\qquad r_{A_u}=r_{A_v}=N-2,\qquad r_S=N-3
\]
である。包含とその右随伴を
$a_R:\mathcal K_R\hookrightarrow\mathcal C_R$、$p_R$ と書く。
以下で使うのは、これらの proper な線束模型、線束上の Serre 双対、
および cubic rotation の指定周期 $\varrho_R^3\simeq[2]$ である。
これらは命題\ref{mci:prop:ambient-line-model}の $m=2,1$ の模型と、
補題\ref{ebk:lem:boundary}の固定境界模型である。
共有中心上の相対 potential 解消と中心方向の Koszul 解消、
命題\ref{mcp:prop:free-quotient}の自由基底と正則性により、
この模型は四つの代数とその制限について両立する。
全 coherent 圏が proper であるという仮定は置かない。

\begin{proposition}[残余圏上の Beck--Chevalley 同型]
\label{tcr:prop:beck-chevalley}
実導来制限は関手
\[
 \begin{array}{ccc}
 \mathcal K_A&\xrightarrow{F_u}&\mathcal K_{A_u}\\
 {\scriptstyle F_v}\downarrow&&\downarrow{\scriptstyle G_v}\\
 \mathcal K_{A_v}&\xrightarrow{G_u}&\mathcal K_S
 \end{array}
\]
を与え、導来 tensor の結合子による指定同型
\begin{equation}\label{tcr:eq:left-square}
                  G_vF_u\simeq G_uF_v=:H
\end{equation}
を持つ。四つの関手は左右の随伴を持つ。
式\eqref{tcr:eq:left-square}の二つの右 mate
\begin{equation}\label{tcr:eq:bc}
       F_vF_u^R\longrightarrow G_u^RG_v,
       \qquad F_uF_v^R\longrightarrow G_v^RG_u
\end{equation}
は同型である。これらは任意に選んだ関手同型ではなく、
実随伴の単位・余単位で定まる射である。
一次の自己交叉が非分裂でも成立する。
\end{proposition}

\begin{proof}
各中心変数は、もう一方を零とした後にも正則である。
線束の二項解消から、制限・inflation とその随伴は全線束生成圏を保存する。
$M\in\mathcal K_A$ の $u$ 制限は $[M(-1)\xrightarrow{u}M]$ で計算できる。
$0\le j\le r_A-2$ に対して、その $\OO(j)$ からの Hom は
$\OO_A(j+1),\OO_A(j)$ から $M$ への Hom であり、共に零である。
従って制限は、さらに projection を施さずに $\mathcal K_{A_u}$ へ入る。
他の三辺も同じである。
二つの順序の制限は $-\otimes_A^{\mathbf L}S$ そのものであり、
式\eqref{tcr:eq:left-square}とその結合 coherence を得る。

$i_u:A_u\to A$、$i_v:A_v\to A$ は inflation の記号を兼ね、
$j_u:S\to A_v$、$j_v:S\to A_u$ とする。
例えば $F_u^R=p_Ai_{u*}$、$G_u^R=p_{A_v}j_{u*}$ である。
$M\in\mathcal K_{A_u}$ に対し、右 projection の三角を
\begin{equation}\label{tcr:eq:mutation-correction}
     a_Ap_Ai_{u*}M\longrightarrow i_{u*}M\longrightarrow E_M
\end{equation}
と書くと、
\begin{equation}\label{tcr:eq:single-correction}
                         E_M\in\langle\OO_A(-1)\rangle
\end{equation}
である。実際、Serre 双対による負の例外列は
$\OO_A(-r_A),\ldots,\OO_A(-1)$ であり、
\[
 R\operatorname{Hom}_A(i_{u*}M,\OO_A(j))
   =R\operatorname{Hom}_{A_u}(M,\OO_{A_u}(j+1))[-1]
\]
は $-r_A\le j\le-2$ で零である。
これは $\mathcal K_{A_u}$ が
$\OO_{A_u}(-(r_A-1)),\ldots,\OO_{A_u}(-1)$ の左直交にあるためである。
従って右 mutation の補正には最後の一線束だけが残る。

正則列の二項解消は Tor-independent な基底変更同型
\begin{equation}\label{tcr:eq:ambient-bc}
                i_v^*i_{u*}M\simeq j_{u*}j_v^*M
\end{equation}
を与える。ここで star は導来制限を表す。
式\eqref{tcr:eq:mutation-correction}へ $p_{A_v}i_v^*$ を施す。
補正は $p_{A_v}\OO_{A_v}(-1)=0$ によって消えるので、
\[
 F_vF_u^RM
   \xrightarrow{\ \sim\ }p_{A_v}i_v^*i_{u*}M
   \xrightarrow{\ \sim\ }p_{A_v}j_{u*}j_v^*M
   =G_u^RG_vM
\]
を得る。最初の射は projection の余単位、次の射は
式\eqref{tcr:eq:ambient-bc}そのものである。
この合成を $G_u\dashv G_u^R$ で transpose すると
\[
 G_uF_vF_u^R\simeq G_vF_uF_u^R
                   \xrightarrow{G_v\epsilon_{F_u}}G_v
\]
となる。従って、これは式\eqref{tcr:eq:bc}の指定 mate である。
$u,v$ を交換すればもう一方も同型になる。
この証明は自己交叉の分裂を使っていない。
\end{proof}

\begin{proposition}[共同 comonad と復元に必要な追加データ]
\label{tcr:prop:joint-comonad}
境界 $\mathcal K_S$ 上の軸 comonad を
\[
             K_u=G_uG_u^R,\qquad K_v=G_vG_v^R
\]
と書く。実随伴正方形は、全 coherence を持つ共同 comonad
$K_{uv}=HH^R$ と、その underlying kernel の同型
\begin{equation}\label{tcr:eq:joint-factorization}
                    K_{uv}\simeq K_vK_u\simeq K_uK_v
\end{equation}
を与える。この同型によって得る交換射と、その余単位・余乗法との
高次互換性を含むデータを、ここでは\emph{coherent な分配データ}と呼ぶ。
それを保持すれば、$K_{uv}$ の coKleisli 圏の thick 冪等完備化から
$\mathcal K_A$ を復元できる。
二つの comonad $K_u,K_v$ の個別の同型だけでは、この追加データは指定されない。
\end{proposition}

\begin{proof}
右随伴の合成と式\eqref{tcr:eq:bc}から
\[
 \begin{aligned}
 HH^R
   &=G_uF_vF_u^RG_v^R
     \simeq G_uG_u^RG_vG_v^R=K_uK_v,\\
 HH^R
   &=G_vF_uF_v^RG_u^R
     \simeq G_vG_v^RG_uG_u^R=K_vK_u
 \end{aligned}
\]
を得る。ここで合成の同型は、左随伴の式\eqref{tcr:eq:left-square}と
その右随伴への mate を同時に使ったものである。
交換射は、これら二つの同型を $HH^R$ を介して合成して得る。
余単位 $HH^R\to\mathbf1$ と余乗法 $HH^R\to HH^RHH^R$ も移す。
例えば $K_{uv}\to K_v$ は
$G_vF_uF_u^RG_v^R\xrightarrow{G_v\epsilon_{F_u}G_v^R}G_vG_v^R$
である。もう一方も同様である。

通常の随伴正方形から分配則と合成 comonad を得る形式論は
\cite[Section~3, Proposition]{CMOBeckDistributiveLaws}にある。
本稿では実 dg 核、導来 tensor の結合子、projection の随伴を固定しているため、
移されるのは homotopy category 上の二つの等式だけではなく、
$HH^R$ が元から持つ全 comonad 構造である。
裸の交換射からその高次構造が一意に定まるとは主張しない。

最後に生成性を確かめる。
命題\ref{cmo:prop:central-linear-spherical}の二項解消・mutation の証明は、
例外線束列の長さを上記の値に替えると四辺に適用できる。
特に各辺 $F$ の単位の fiber は $\varrho^{-1}$、その三乗は $[-2]$ である。
proper 性により $M[-2q]\to M$ は $q\gg0$ で零になり、
反復した単位三角から
$\mathcal K_{\mathrm{source}}=\operatorname{thick}(F^R\mathcal K_{\mathrm{target}})$
を得る。二回適用して
\[
                  \mathcal K_A=\operatorname{thick}(H^R\mathcal K_S)
\]
となる。境界の対象 $E,F$ について
\[
 R\operatorname{Hom}_{\mathcal K_A}(H^RE,H^RF)
    =R\operatorname{Hom}_{\mathcal K_S}(K_{uv}E,F)
\]
であり、積は実余乗法で与えられる。これが主張した coKleisli 復元である。
全 coherent 圏への拡張や、源の smoothness はこの復元の仮定ではない。
\end{proof}

\begin{lemma}[二つの法線の Koszul 行列因子分解に現れる混合積]
\label{tcr:lem:mixed-normal-products}
$R$ を可換次数付き $\C$-代数とし、$A_2,B_2,C_2\in R_1$、
$d,e,f,h\in R_0$ とする。
\begin{equation}\label{tcr:eq:normal-potential}
 q=u^2A_2+uvB_2+v^2C_2+du^3+eu^2v+fuv^2+hv^3
\end{equation}
に対して、$q$ の Koszul 行列因子分解の dg End は
$R[t,t^{-1}]$ 上自由階数四の cohomology を持つ。
ここで $|t|=(2,-3)$ は cohomological/internal 双次数である。
二つの閉元の class を $\theta_u,\theta_v$、$|\theta_i|=(1,-1)$ とすると、
二項積は
\begin{equation}\label{tcr:eq:normal-clifford}
 \theta_u^2=-tA_2,\qquad
 \theta_u\theta_v+\theta_v\theta_u=-tB_2,\qquad
 \theta_v^2=-tC_2
\end{equation}
となる。以下の Koszul contraction から得る minimal 模型では、
$\theta=r\theta_u+s\theta_v$ に対し、共通の $A_\infty$ 符号規約を除いて
\begin{equation}\label{tcr:eq:normal-cubic}
       m_3(\theta,\theta,\theta)
              =\pm t(dr^3+er^2s+frs^2+hs^3)
\end{equation}
である。従って、軸への制限で消える $uv$ と $u^2v,uv^2$ の係数は、
この指定した実 dg 模型の二項積と三項積にそれぞれ現れる。
\end{lemma}

\begin{proof}
外積演算子 $\varepsilon_u,\varepsilon_v$ と縮約演算子 $\iota_u,\iota_v$ を
二変数外積加群上に取り、
$|\varepsilon_i|=(-1,1)$、$|\iota_i|=(1,-1)$ とする。
それぞれの反交換子は零、$\{\iota_i,\varepsilon_j\}=\delta_{ij}$ である。
対称行列
\[
 (h_{ij})=
 \begin{pmatrix}
 A_2+du+ev/3&B_2/2+eu/3+fv/3\\
 B_2/2+eu/3+fv/3&C_2+fu/3+hv
 \end{pmatrix}
\]
を置き、$a=uh_{uu}+vh_{uv}$、$b=uh_{uv}+vh_{vv}$ とする。
\[
       D=u\iota_u+v\iota_v+t(a\varepsilon_u+b\varepsilon_v)
\]
は次数 $(1,0)$ で $D^2=tq$ を満たす。
従って $[D,-]$ はその End 上の微分になる。
\[
             \Theta_i=\iota_i-t\sum_jh_{ij}\varepsilon_j
\]
と置くと、直接計算で
\[
 [D,\varepsilon_u]=u,\quad [D,\varepsilon_v]=v,
 \quad [D,\Theta_i]=0,
 \quad \Theta_i^2=-th_{ii},\quad
 \Theta_u\Theta_v+\Theta_v\Theta_u=-2th_{uv}
\]
を得る。角括弧は次数付き交換子である。
End の加群基底を $\varepsilon_u,\varepsilon_v,\Theta_u,\Theta_v$ の
順序付き積へ替えると、その複体は $u,v$ の Koszul 複体と
$1,\Theta_u,\Theta_v,\Theta_u\Theta_v$ の自由加群の tensor になる。
多項式環上で $u,v$ は $R$ によらず正則列なので、
cohomology の自由性と式\eqref{tcr:eq:normal-clifford}が従う。

次数を保つ通常の Euler contraction をこの Koszul 複体に取る。
特に $\gamma u+\eta v$ の contraction は
$\gamma\varepsilon_u+\eta\varepsilon_v$ とする。
$\Theta=r\Theta_u+s\Theta_v$ と置き、
\[
 \begin{aligned}
 T_{r,s}=-t\bigl(&
 (dr^2+2ers/3+fs^2/3)\varepsilon_u\\
 &+(er^2/3+2frs/3+hs^2)\varepsilon_v\bigr)
 \end{aligned}
\]
とすると、
\[
 \begin{aligned}
 [D,T_{r,s}]&=\Theta^2+t(A_2r^2+B_2rs+C_2s^2),\\
 \Theta T_{r,s}+T_{r,s}\Theta
       &=-t(dr^3+er^2s+frs^2+hs^3).
 \end{aligned}
\]
第一式は二項積の代表とその cohomology class の差の homotopy、
第二式はその三項移送そのものである。
strict unit の contraction は零に取れるため、二項積の scalar 部分は
この三項式へ追加項を与えない。
これで式\eqref{tcr:eq:normal-cubic}を得る。
\end{proof}

\begin{remark}[有限計算と全分配データの区別]
\label{tcr:rem:balancing-scope}
補題\ref{tcr:lem:mixed-normal-products}は、例えば可換式 $W=G+q$ における
法線方向の実 Koszul 行列因子分解を計算している。
全 active 方向を含む残余圏の任意の gauge を割った moduli で、
表示した混合係数が全て独立であることまで証明したものではない。
また、$uQ_u+vQ_v$ を持つ非分裂の一次自己交叉をこの計算から除去することもしない。
命題\ref{tcr:prop:beck-chevalley}と\ref{tcr:prop:joint-comonad}は
そのような非分裂の場合にも成立する。

二重 Koszul 解消が加群として有限自由であることから、
共同 comonad の比較に必要な balancing tower が有限表示であるとは結論できない。
既に一変数の相対対角 $z=u_L-u_R$ の associative cofibrant 解消には、
\[
 |q_m|=(1-2m,m),\qquad dq_1=z,\qquad
 dq_m=\sum_{i+j=m}q_iq_j\quad(m\ge2)
\]
という無限列が現れる。
内部次数と三次周期から高 arity の消滅を評価するには、
その評価対象が実際の coherent 分配データの tower であるという模型同定が別に必要である。
本小節はその同定や全高次 fiber の可縮性を主張しない。

\path{computations/verify_two_center_restrictions.py}は、
$4\times4$ 行列による上の法線恒等式と、次の例の二変数部分を六項目で再現する。
残余 mutation、Beck--Chevalley coherence、smoothness の主張は本文の証明による。
\end{remark}

\begin{example}[同一の滑らかな軸と異なる全残余圏]
\label{tcr:ex:same-axes}
$G(x_1,\ldots,x_N)$ を滑らかな射影三次の方程式とし、
\[
 \begin{aligned}
 W_0&=G+u^3+v^3,\\
 W_1&=G+u^3+v^3-u^2v-uv^2
      =G+(u+v)(u-v)^2
 \end{aligned}
\]
と置く。$X_i=V(W_i)\subset\mathbb P^{N+1}$ とすると、
両者の $u=0$、$v=0$ の軸超平面切断と、それらから
$Z=V(G)$ への実制限随伴は同一である。
二つの軸切断はともに滑らかである。
従って境界上の二つの軸 spherical comonad は、各々の全構造を含めて同一に取れる。
しかし $\mathcal K_{X_0}$ は smooth で、残余線束圏 $\mathcal K_{X_1}$ は smooth でない。
特に個別の軸 comonad の同定から、共同 comonad の coherent 分配データは決まらない。
\end{example}

\begin{proof}
斉次性と Euler 公式から、$\partial G=0$ の affine 解は原点だけである。
$W_0$ の全偏微分は $\partial G,3u^2,3v^2$ なので、
それらは射影空間で同時に消えず、$X_0$ は滑らかである。
同じ計算で $V(G+u^3)$、$V(G+v^3)$ も滑らかである。
一方、
\[
 \partial_u\bigl((u+v)(u-v)^2\bigr)=(u-v)(3u+v),\qquad
 \partial_v\bigl((u+v)(u-v)^2\bigr)=-(u-v)(u+3v)
\]
である。従って $p=[0:\cdots:0:1:1]\in X_1$ では全射影偏微分が消え、
$p$ は特異点である。$v=1$、$u=1+s$ の局所方程式は
\[
                          G(\mathbf x)+(2+s)s^2=0.
\]
これは極大イデアルの二乗に入り、局所環の embedding dimension は $N+1$、
Krull dimension は $N$ なので、局所環は正則でない。

$u=0$ または $v=0$ とすると二つの式の混合項は全て消える。
従って軸の次数付き代数だけでなく、$S=\C[\mathbf x]/(G)$ への二つの商写像も同じである。
線束列、projection、導来制限、inflation を同一の dg 模型で取れば、
実随伴の単位・余単位と、それらから作る軸 comonad も全て同一になる。

最後に残余線束圏の smoothness を比較する。
射影 hypersurface では全線束 twist は $\operatorname{Perf}(X_i)$ を生成し、
標準例外列は
\[
 \operatorname{Perf}(X_i)
   =\langle\mathcal K_{X_i},\OO_{X_i},\ldots,\OO_{X_i}(N-2)\rangle
\]
を与える。$X_0$ は滑らかなので、その admissible 成分
$\mathcal K_{X_0}$ は smooth である。
もし $\mathcal K_{X_1}$ が smooth なら、この proper な成分と有限例外列の
gluing 双加群は perfect であり、$\operatorname{Perf}(X_1)$ も smooth になる。
ここで proper な加群が smooth な圏上で perfect となることは、
有限な対角解消をその加群と tensor すれば分かる。
有限例外列との gluing にも同じ対角解消を用いる。

この帰結は点 $p$ で矛盾する。
実際、$\operatorname{Perf}(X_1)$ の有限生成対象 $T$ に対し、
$R\operatorname{Hom}(T,\OO_p)$ は $\C$ 上 perfect である。
$\operatorname{Perf}(X_1)$ が smooth なら、この加群は
$R\operatorname{End}(T)$ 上も perfect となり、$\OO_p$ は perfect complex で表される。
全線束 twist の生成性により、その表現は実際に $\OO_p$ と同型である。
しかし局所環が正則でないため、その residue field は有限 projective dimension を持たない。
これは $\OO_p$ の perfect 性に反する。
よって $\mathcal K_{X_1}$ は smooth でない。
この議論は特異な全 coherent 残余圏の比較を必要としない。
\end{proof}

